% !TEX program = xelatex
% 用户指令语用特征对任务型交互智能体执行效果的影响
% 基于 AgentNet 大规模 GUI 交互轨迹数据的实证研究三部曲

\documentclass[12pt,a4paper,openany,oneside]{ctexrep}

% ── Font & Encoding ──
\setCJKmainfont{SimSun}
\setCJKsansfont{SimHei}
\setCJKmonofont{KaiTi}
\setmainfont{Times New Roman}
\setsansfont{Arial}
\setmonofont{Courier New}

% ── Page Layout ──
\usepackage[paper=a4paper,
            top=2.54cm,
            bottom=2.54cm,
            left=3.18cm,
            right=3.18cm]{geometry}

% ── Spacing ──
\usepackage{setspace}
\onehalfspacing
\setlength{\parindent}{2em}

% ── Headers & Footers ──
\usepackage{fancyhdr}
\setlength{\headheight}{27pt}
\addtolength{\topmargin}{-12pt}
\pagestyle{fancy}
\fancyhf{}
\fancyhead[C]{\small\kaishu\leftmark}
\fancyfoot[C]{\thepage}
\renewcommand{\headrulewidth}{0.4pt}

% ── Tables ──
\usepackage{booktabs}
\usepackage{longtable}
\usepackage{threeparttable}
\usepackage{adjustbox}
\usepackage{array}
\newcommand{\tabindent}{0.5em}

% ── Graphics ──
\usepackage{graphicx}
\usepackage{float}

% ── Math ──
\usepackage{amsmath}
\usepackage{amssymb}

% ── Captions ──
\usepackage{caption}
\captionsetup[figure]{
  font=small,
  labelfont=bf,
  justification=centering,
  skip=6pt
}
\captionsetup[table]{
  font=small,
  labelfont=bf,
  justification=centering,
  skip=6pt,
  position=above
}

% ── Hyperlinks ──
\usepackage[colorlinks=true,
            linkcolor=blue,
            citecolor=blue,
            urlcolor=blue]{hyperref}

% ── Lists ──
\usepackage{enumitem}
\setlist{nosep, leftmargin=2em, itemsep=2pt, parsep=2pt}

% ── TOC ──
\usepackage{tocbibind}
\setcounter{tocdepth}{3}

% ── Miscellaneous ──
\usepackage{xcolor}
\usepackage{csquotes}

% ── Document Info ──
\title{用户指令语用特征对任务型交互智能体执行效果的影响}
\author{}
\date{\\today}

\begin{document}

% ────────────────────────────────────
%  Cover / Title Page
% ────────────────────────────────────

\begin{titlepage}
\thispagestyle{empty}

\begin{center}

\vspace*{4cm}

{\zihao{0}\heiti 用户指令语用特征对\\[6pt]
任务型交互智能体执行效果的影响}

\vspace{20pt}

{\zihao{3}\kaishu —— 基于 AgentNet 大规模 GUI 交互轨迹数据的\\[4pt]
实证研究三部曲}

\vspace{40pt}

{\zihao{4}}
\vfill

{\zihao{4}\today}

\end{center}

\end{titlepage}

% ── TOC ──
\chapter*{综合摘要}
\tableofcontents
\newpage




\newpage

\newpage

\chapter{第1部分　完成标准作为终止判准：用户指令验证显性度与任务型交互智能体执行效率——基于 AgentNet 交互轨迹的实证研究}


\subsection{1 引言}


\subsubsection{1.1 研究背景}


以大语言模型（Large Language Model, LLM）为核心、融合视觉感知的任务型交互智能体（task-oriented interactive agent），近年来在人机交互、自动化办公、软件测试与数字助手等领域的应用快速扩展。典型系统如 Mind2Web（Yang et al., 2022）、WebGPT（Pal et al., 2022）以及 AUSTIN（Yuan et al., 2023），展示了智能体在网页浏览、文件管理与复杂多步骤任务中的执行能力；AgentNet（Zhang et al., 2024）等大规模真实交互轨迹数据集的发布，则为从自然语言指令维度理解人机协同效能提供了宝贵的经验基础。

从过程视角审视，任务型交互智能体的执行可抽象为一条通用链：\textbf{用户任务表达 → 意图解析 → 动作规划 → 环境交互 → 状态评估 → 执行终止}。在这一链条中，终止环节决定智能体何时停止当前动作序列并返回结果，其重要性不亚于规划与执行本身。然而，该环节在既有研究中受到的关注远不及动作规划与任务分解。

\subsubsection{1.2 问题提出：终止不确定性}


与自然语言处理中常见的"生成到句末即停"不同，任务型智能体的终止必须依赖对外部任务状态是否达标的判断。当用户指令仅给出高层目标——如"整理桌面"——而未说明完成判据——如"所有截图文件移入 Screenshots 文件夹且桌面剩余文件不超过 5 个"——智能体便面临一种结构性的\textbf{终止不确定性}（termination uncertainty）。这种不确定性可能表现为两种对立的偏差：

\begin{enumerate}
\item[1.] \textbf{提前终止}（premature termination）：智能体在未完全达成潜在目标时即停止执行，导致任务失败；
\item[2.] \textbf{过冲探索}（overshooting）：智能体因缺乏明确判据而反复验证、冗余搜索，超出必要执行步数。
\end{enumerate}


上述终止偏差的存在，将研究焦点指向一个根本性问题：用户指令中的何种语言信息能够为智能体提供可靠的停止信号？

\subsubsection{1.3 核心研究问题与论点}


本研究以 AgentNet GUI 交互轨迹为经验边界，将用户指令中的完成标准/验证条件重新概念化为智能体的\textbf{"终止判准"}（stopping cue）。核心研究问题为：当指令明确包含"检查/确认/确保……"等验证性表述或"完成/就绪/正确"等目标状态时，智能体是否更善于将自然语言目标转译为可操作的停止条件，并与更高的执行效率相关联？

本文的基本论点是：\textbf{任务型交互智能体必须将自然语言目标转译为可执行的停止条件（executable stopping condition）；明确的验证标准与完成状态能够降低智能体对"是否已达目标"的不确定性，并与更少的冗余探索和提前终止相关联。} 需要强调的是，本文基于观察性交互轨迹数据，所有发现均描述统计关联，而非因果效应。

\subsection{2 理论基础与研究假设}


本研究的理论框架由三条相互关联的理论线索构成：言语行为理论提供指令语力的分析工具，关联理论解释信息加工的认知经济性，行动终止理论则为停止规则的操作化提供理论依据。三条线索共同指向一个核心推论：完成标准作为指令中的显性验证信息，构成智能体内部终止规则的语言基础。

\subsubsection{2.1 言语行为理论与指令中的验证行为}


言语行为理论（Speech Act Theory）的核心洞见在于：语言不仅是描述世界的工具，也是执行行为的手段（Austin, 1962; Searle, 1969）。在任务型交互情境中，用户指令本质上是一种\textbf{指令式言语行为}（directive speech act），其语力（illocutionary force）在于引导智能体执行特定操作序列。

对该理论框架的进一步审视揭示了一个重要但常被忽视的维度：指令内部承载着不同层次的语力信息。当指令包含"检查""确认""确保"等动词时，用户不仅指示智能体完成操作动作，还同时规定了验证行为——要求智能体将当前状态与目标状态进行比较并做出判断。当指令出现"完成""正确""就绪"等状态描述时，则为智能体提供了终止执行所需的语义目标。换言之，完成标准与验证条件并非指令的附属修饰，而是指令语力的构成性组成部分，直接影响智能体对任务边界的识别与终止时机的判断。

\subsubsection{2.2 关联理论与认知效应优化}


关联理论（Relevance Theory）强调，交际双方在信息处理中追求以最小认知努力获取最大认知效应（Sperber \& Wilson, 1995）。这一"认知经济性"原则在任务型交互智能体的语境中具有直接的解释力。

具体而言，明确的完成标准能够同时作用于认知效应的两端：一方面，它降低了意图识别与状态评估的认知成本——智能体无需从模糊目标中推断停止条件；另一方面，它提高了智能体对任务完成状态的判断精度——验证动作的执行直接产生可观测的状态证据。二者共同预期：含完成标准的指令与更少的执行步数和更高的完成率相关联。

然而，关联理论同时内含一个值得警惕的逻辑：根据关联原则，过高的验证密度可能在提供冗余约束的同时限制智能体的灵活推理空间，从而可能削弱甚至逆转验证信息的正向效应。这提示验证密度与任务结果之间可能存在非线性关系，构成一个有待经验检验的开放问题。

\subsubsection{2.3 行动终止理论与停止规则}


行动终止/停止规则理论关注行为主体在何种条件下停止当前动作并转入下一阶段（Gigerenzer et al., 1999; Payne et al., 1993）。该理论区分两类终止规则：\textbf{内部终止}（达到预设目标状态时停止）与\textbf{外部终止}（用户干预或超时等外部信号触发停止）。

在任务型交互智能体的执行情境中，当指令明确提供完成标准时，智能体可依据内部终止规则自主判断停止时机，减少对外部启发式规则的依赖，从而抑制不必要的后续步骤。反之，当指令缺乏完成标准时，模糊的终止条件迫使智能体依赖启发式停止规则——例如"操作步骤达到一定数量"或"连续若干步无状态变化"——这类规则缺乏对任务目标状态的直接参照，容易导致过冲或提前终止。由此，完成标准构成了内部终止规则的关键语言线索。

\subsubsection{2.4 研究假设}


基于上述三条理论线索的交叉论证，本研究提出以下假设：

\textbf{H1a（完成率假设）}：包含完成标准/验证条件的任务，其完成率显著高于不包含完成标准的任务。

\textbf{H1b（执行效率假设）}：包含完成标准的任务执行步数显著少于不包含完成标准的任务。

\textbf{H1c（终止过程假设）}：包含完成标准的任务在终止阶段具有更低的冗余比例与错误比例。

上述 H1a–H1c 构成本研究的核心假设，直接检验"完成标准作为终止判准"这一理论框架的基本预测。

\textbf{H2（验证类型差异，探索性）}：不同验证语言形式（动词指令式、状态目标式、条件验证式、混合式）对执行终止过程的影响可能存在差异。该假设为探索性，因类型分布严重失衡。

\textbf{H3（密度非线性，补充性探索）}：验证关键词密度与任务完成率可能呈非线性关系。该假设源于关联理论中认知成本与效应之权衡的逻辑推论，但在当前数据情境下作为补充性分析。

\textbf{H4（位置效应，探索性）}：验证表述位于指令不同位置时，其终止引导效应可能存在差异。该假设基于信息结构理论，但 AgentNet 数据中验证表述的分布特征限制了其检验力。

\subsection{3 研究方法}


\subsubsection{3.1 数据来源与样本口径}


本研究使用公开的 AgentNet 数据集（Zhang et al., 2024），原始数据包含 $N=22{,}811$ 条自然语言 GUI 任务及其交互轨迹。该数据集为真实人机/人–智能体交互的观察性记录，非实验操纵数据。任务完成状态、步骤正确性、冗余性等字段由数据集原始标注或后续规则化处理得到。

为澄清研究结论的适用范围，本研究将分析样本明确区分为四个层次：

\begin{enumerate}
\item[1.] \textbf{原始全量样本}（$N=22{,}811$）：用于描述总体分布与完成标准二值特征的标注；
\item[2.] \textbf{含完成标准原始子集}（$n=2{,}869$，即 `has\_completion\_criteria == 1`）：用于验证语言形式的分类描述；
\item[3.] \textbf{回归分析样本}（$N=20{,}484$）：在剔除 `task\_completed` 与 `task\_difficulty` 缺失值后构成，其中含完成标准任务 $n=2{,}684$，不含完成标准任务 $n=17{,}800$；
\item[4.] \textbf{位置分析样本}（$N=2{,}512$）：含完成标准且能定位到验证关键词的子集。
\end{enumerate}


\subsubsection{3.2 变量操作化}


\paragraph{3.2.1 核心解释变量}


核心解释变量为二值变量 `has\_completion\_criteria`（是否包含完成标准/验证条件），该变量由数据集原始标注或前期规则化提取得到，本质上是一种基于关键词与句法启发式的规则化特征。需特别说明的是，规则化特征可能存在误报或漏报——例如，某些含完成标准的表述可能未命中关键词，而某些关键词在不同语境下未必表达完成判据；此外，原始标注可能偏向"检查/确认"类动词，导致状态目标式与条件验证式样本稀少。因此，本文将 `has\_completion\_criteria` 视为一种可观测的语言线索，而非对"指令质量"的全面度量。

对 `has\_completion\_criteria == 1` 的指令，依据关键词与句法窗口规则进一步细分为五类：

\begin{itemize}
\item \textbf{V1 动词指令式}：出现 `check` / `verify` / `confirm` / `validate` / `ensure` / `make sure` 等动词，直接指示验证动作；
\item \textbf{V2 状态目标式}：出现 `done` / `complete` / `correct` / `proper` / `finished` / `ready` 等状态词，描述完成后的目标状态；
\item \textbf{V3 条件验证式}：出现 `until` / `when` / `once` / `after` 等条件词，其中 `until` 直接视为完成条件，`when` / `once` / `after` 仅在周围 15 词窗口内同时出现 V1/V2 关键词时才视为完成条件；
\item \textbf{V4 混合式}：同时命中 V1/V2/V3 中至少两类；
\item \textbf{V0 不可分类/伪阳性}：`has\_completion\_criteria == 1` 但未命中任何 V1/V2/V3 关键词。
\end{itemize}


验证密度 `verification\_density` 定义为验证关键词出现次数除以指令总词数。验证位置 `verification\_position\_region` 按字符相对位置划分为开头（< 0.25）、中间（0.25–0.75）与末尾（≥ 0.75）三类。

\paragraph{3.2.2 因变量}


\begin{itemize}
\item `task\_completed`：任务完成状态（0/1）；
\item `total\_steps`：总执行步数（计数变量）；
\item 终止阶段指标（见 3.3 节）。
\end{itemize}


\paragraph{3.2.3 控制变量}


\begin{itemize}
\item `task\_difficulty`：任务难度（1–10 连续评分）；
\item `word\_count`：指令词数；
\item `os`：操作系统（ubuntu / win\_mac）。
\end{itemize}


\subsubsection{3.3 终止阶段过程指标的构建}


为从过程层面刻画完成标准对终止行为的影响，本研究基于 `data/processed/analysis\_data.parquet` 与 `data/interim/steps.parquet` 计算以下终止阶段指标：

\begin{enumerate}
\item[1.] \textbf{最后 25\% 步冗余比例}（`last25\_redundant\_prop`）：对每条任务，取步数最后 25\%（向上取整，至少 1 步）的步骤，计算其中 `last\_step\_redundant == True` 的比例；
\item[2.] \textbf{最后 25\% 步错误比例}（`last25\_incorrect\_prop`）：同上，但计算 `last\_step\_correct == False` 的比例；
\item[3.] \textbf{最后一步是否冗余}（`final\_step\_redundant`）：任务最后一步 `last\_step\_redundant == True`；
\item[4.] \textbf{最后一步是否错误}（`final\_step\_incorrect`）：任务最后一步 `last\_step\_correct == False`；
\item[5.] \textbf{最后 25\% 步是否出现冗余/错误}（`last25\_any\_redundant`、`last25\_any\_incorrect`）：二值指标；
\item[6.] \textbf{异常短轨迹}（`abnormal\_short`）：针对未完成任务，若其总步数不高于已完成任务总步数的第 5 百分位数（本数据中为 9 步），则标记为异常短轨迹，作为提前终止失败的描述性指标。
\end{enumerate}


上述指标的计算结果保存于 `revised/paper01\_tables/terminal\_phase\_metrics\_by\_task.csv`，按 `has\_completion\_criteria` 分组汇总于 `revised/paper01\_tables/terminal\_phase\_summary\_by\_criteria.csv`。

\subsubsection{3.4 分析模型}


针对不同因变量的分布特征与假设结构，本研究采用以下分析策略：

\begin{itemize}
\item \textbf{任务完成率}：采用 Logistic 回归，报告比值比（OR）及其 95\% Wald 置信区间；
\item \textbf{执行步数}：采用负二项回归（NB2），报告发生率比（IRR）及 95\% CI，并报告离散参数 alpha 以评估过度离散程度；
\item \textbf{验证密度的非线性效应}：在 Logistic 模型中加入 `verification\_density` 与 `verification\_density\_sq`，采用 AIC 与似然比检验（LRT）比较线性与二次模型。密度分析为补充性探索，结果置于 4.5 节；
\item \textbf{位置效应}：采用卡方检验与 Bonferroni 校正后的成对比例 Z 检验；
\item \textbf{终止阶段指标}：采用 Welch t 检验（连续比例）与双比例 Z 检验（二值指标）进行组间比较。
\end{itemize}


所有分析在 Python 中完成，显著性水平 α = 0.05。鉴于数据的观察性性质，所有回归系数均解释为条件关联，而非因果效应。

\subsection{4 结果}


\subsubsection{4.1 描述性统计与验证类型分布}


在回归样本（$N=20{,}484$）中，总体任务完成率为 79.22\%。无完成标准任务 $n=17{,}800$，完成率为 78.76\%，平均执行步数为 20.61（SD = 17.12）；有完成标准任务 $n=2{,}684$，完成率为 82.27\%，平均执行步数为 19.64（SD = 11.32）。初步描述性统计表明，含完成标准的任务在完成率上具有约 3.5 个百分点的优势，且平均执行步数更低——这一原始差异为后续回归分析提供了描述层面的初步证据，但其是否独立于任务难度、指令长度等混淆变量，需通过模型检验加以确认。

\textbf{表 1} 报告了验证语言形式在含完成标准原始子集（$N=2{,}869$）中的分布。一个显著特征是类型的极端失衡：动词指令式（V1）占绝对主导（82.9\%），混合式（V4）次之（9.7\%），不可分类/伪阳性（V0）占 7.1\%，而状态目标式（V2）与条件验证式（V3）分别仅占 0.2\% 与 0.0\%。这一严重失衡的分布意味着 H2 的组间推断力极为有限，验证类型差异仅作描述性展示，不宜进行统计显著性比较。


\begin{table}[htbp]
\centering
\caption{表 1：验证语言形式在含完成标准任务中的分布（$N=2{,}869$）}
\small
\adjustbox{max width=\textwidth}{
\begin{tabular}{l l r r l}
\toprule
类型 & 标签 & $n$ & 比例 & 典型样例（节选） \\
\midrule
V0 & 不可分类/伪阳性 & 205 & 7.1\% & "…proceed to checkout" / "Export and copy the final content." \\
V1 & 动词指令式 & 2,379 & 82.9\% & "…finally look at the file content and make sure the file exists" \\
V2 & 状态目标式 & 5 & 0.2\% & "…set confirmation completed" \\
V3 & 条件验证式 & 1 & 0.0\% & "…trace the edges of the egg in the picture until they close…" \\
V4 & 混合式 & 279 & 9.7\% & "…After the installation is completed, restart the application." \\
\bottomrule
\end{tabular}
}
\end{table}

上述分布特征揭示了一个值得注意的模式：在真实用户指令中，完成标准的表达方式高度集中在直接动词指令上，状态描述和条件表达极为罕见。这既可能反映用户的自然表达倾向，也可能反映规则化标注对特定关键词（尤其是 V1 类动词）的敏感性。

\begin{figure}[htbp]
\centering
\includegraphics[width=0.92\textwidth]{paper01_figures/fig1_验证类型分布.png}
\caption{图 1　含完成标准任务中的验证类型分布}
\end{figure}


\textbf{图 1} 以可视化方式呈现了验证类型的分布特征，直观展示了 V1 的主导地位与其他类型的稀少性。

\subsubsection{4.2 完成标准对任务完成率的影响}


\textbf{表 2} 报告了 Logistic 回归结果（$N=20{,}484$）。在控制任务难度、指令词数与操作系统后，`has\_completion\_criteria` 的 OR 为 1.255（95\% CI [1.127, 1.396]，$p<.001$），表明含完成标准的任务完成 odds 提高约 25.5\%。该结果支持 H1a，为"完成标准作为终止判准"提供了统计层面上的证据。

操作系统方面，win\_mac 平台的完成率显著低于 Ubuntu 参考组（OR = 0.554，$p<.001$）；任务难度的效应不显著（OR = 0.983，$p=.206$）；指令词数每增加一词，完成 odds 提高约 0.2\%（OR = 1.002，$p=.026$）——这一微弱但显著的效应提示，较长的指令可能包含了更多任务相关信息，从而略微提升完成概率。


\begin{table}[htbp]
\centering
\caption{表 2：任务完成率影响因素的 Logistic 回归结果（$N=20{,}484$）}
\small
\adjustbox{max width=\textwidth}{
\begin{tabular}{l r r}
\toprule
预测变量 & OR [95\% CI] & $p$ \\
\midrule
Intercept & 6.413 [5.421, 7.586] & <.001 \\
`has\_completion\_criteria` & 1.255 [1.127, 1.396] & <.001 \\
`task\_difficulty` & 0.983 [0.958, 1.009] & .206 \\
`word\_count` & 1.002 [1.000, 1.004] & .026 \\
`C(os)[T.win\_mac]` & 0.554 [0.492, 0.623] & <.001 \\
\bottomrule
\end{tabular}
}
\end{table}

\begin{figure}[htbp]
\centering
\includegraphics[width=0.92\textwidth]{paper01_figures/fig2_完成率与步数对比.png}
\caption{图 2　是否包含完成标准的任务完成率与执行步数对比}
\end{figure}


\textbf{图 2} 以双面板呈现了两组任务在完成率和执行步数上的对比，直观显示含完成标准任务在两项指标上的优势方向。

\begin{figure}[htbp]
\centering
\includegraphics[width=0.92\textwidth]{paper01_figures/fig5_逻辑回归系数森林图.png}
\caption{图 5　Logistic 回归系数森林图}
\end{figure}


\textbf{图 5} 以森林图形式呈现了 Logistic 模型中各预测变量的 OR 及 95\% CI，直观展示了 `has\_completion\_criteria` 效应的方向与精度。

\subsubsection{4.3 完成标准对执行成本的影响}


\textbf{表 3} 为负二项回归结果（$N=20{,}484$）。`has\_completion\_criteria` 的 IRR 为 0.883（95\% CI [0.866, 0.900]，$p<.001$），表明在控制其他变量后，含完成标准的任务预期执行步数低约 11.7\%，支持 H1b。

任务难度的 IRR 为 1.211（$p<.001$），即难度每增加 1 分，预期步数增加约 21.1\%，符合"更困难的任务需要更多步骤"的直觉预期。win\_mac 平台的步数略高于 Ubuntu（IRR = 1.079，$p<.001$）。离散参数 α = 0.171，提示步数数据存在适度过度离散，负二项模型较 Poisson 模型更为合适。


\begin{table}[htbp]
\centering
\caption{表 3：执行步数影响因素的负二项回归结果（NB2，$N=20{,}484$）}
\small
\adjustbox{max width=\textwidth}{
\begin{tabular}{l r r}
\toprule
预测变量 & IRR [95\% CI] & $p$ \\
\midrule
Intercept & 5.838 [5.661, 6.020] & <.001 \\
`has\_completion\_criteria` & 0.883 [0.866, 0.900] & <.001 \\
`task\_difficulty` & 1.211 [1.205, 1.217] & <.001 \\
`word\_count` & 1.001 [1.001, 1.002] & <.001 \\
`C(os)[T.win\_mac]` & 1.079 [1.057, 1.100] & <.001 \\
alpha & 1.187 [1.182, 1.192] & <.001 \\
\bottomrule
\end{tabular}
}
\end{table}

\begin{figure}[htbp]
\centering
\includegraphics[width=0.92\textwidth]{paper01_figures/fig6_负二项回归系数森林图.png}
\caption{图 6　负二项回归系数森林图}
\end{figure}


\textbf{图 6} 以森林图呈现了负二项回归中各变量的 IRR 及 95\% CI，直观展示了完成标准与步数减少的关联。

\subsubsection{4.4 终止阶段过程指标}


以上分析从结果层面（完成率）和总量层面（总步数）检验了完成标准的效应。然而，本研究的理论框架预测，完成标准的影响在终止阶段应尤为显著——当智能体接近目标时，明确的验证信息应引导其做出更精准的停止决策。为从过程层面检验这一预测，本研究进一步比较了两组任务在终止阶段的行为差异。

\textbf{表 4} 报告了按 `has\_completion\_criteria` 分组的终止阶段指标描述与组间比较。


\begin{table}[htbp]
\centering
\caption{表 4：终止阶段指标按是否含完成标准的比较}
\small
\adjustbox{max width=\textwidth}{
\begin{tabular}{l r r r r}
\toprule
指标 & 无完成标准（$n=19{,}942$） & 有完成标准（$n=2{,}869$） & 统计量 & $p$ \\
\midrule
最后 25\% 步冗余比例（均值） & 0.0637 & 0.0555 & $t=-3.496$ & <.001 \\
最后 25\% 步错误比例（均值） & 0.1340 & 0.1218 & $t=-3.384$ & <.001 \\
最后 25\% 步出现冗余（比例） & 26.31\% & 23.14\% & $Z=-3.620$ & <.001 \\
最后 25\% 步出现错误（比例） & 44.77\% & 41.65\% & $Z=-3.143$ & .002 \\
最后一步冗余率 & 0.02\% & 0.14\% & $Z=3.193$ & .001 \\
最后一步错误率 & 23.65\% & 23.74\% & $Z=0.104$ & .917 \\
异常短轨迹率（未完成任务） & 2.47\% & 1.19\% & $Z=-4.278$ & <.001 \\
平均总步数 & 19.99 & 19.24 & — & — \\
平均冗余比例 & 0.0728 & 0.0661 & $t=-4.207$ & <.001 \\
平均错误比例 & 0.1188 & 0.1079 & $t=-5.035$ & <.001 \\
\bottomrule
\end{tabular}
}
\end{table}

结果显示，含完成标准任务在最后 25\% 步的冗余比例、错误比例、出现冗余/错误的概率以及整体冗余/错误比例均显著更低；异常短轨迹率（未完成任务中总步数 ≤ 9 步）亦显著更低——后者是 H1c 的关键证据之一，提示缺乏完成标准的指令可能与更频繁的提前终止失败相关联。

最后一步错误率在两组间差异不显著（$p=.917$）；最后一步冗余率虽有统计显著性，但绝对发生率极低（0.02\% vs. 0.14\%），实际意义有限。综合来看，H1c 获得部分支持：完成标准与更低的终止阶段冗余和错误相关联，但这种关联不应被解释为完成标准"导致"了更优的终止行为，而应理解为统计上的共变关系。

\begin{figure}[htbp]
\centering
\includegraphics[width=0.92\textwidth]{paper01_figures/fig_terminal_phase_comparison.png}
\caption{图 7　终止阶段指标按是否含完成标准的比较}
\end{figure}


\textbf{图 7} 以多面板对比了含/不含完成标准任务在终止阶段各项指标上的差异，直观呈现了完成标准与更优终止行为的关联模式。

\subsubsection{4.5 验证密度：补充性分析}


为检验 H3，本研究比较了线性与二次 Logistic 模型。线性模型中，`verification\_density` 的系数显著为正（OR = 25.619，$p<.001$），即验证密度越高，完成率越高。二次模型中，二次项系数不显著（OR = $2.66\times10^{-5}$，$p=.654$），似然比检验 χ² = 0.197，$p=.657$，AIC 亦更支持线性模型（线性 AIC = 20811.69，二次 AIC = 20813.49，ΔAIC = −1.80）。

据此，H3 未获支持：在当前操作化下，未发现验证密度与完成率之间存在稳健的倒 U 形关系。这一结果与关联理论的预测方向不完全一致——验证信息的增加并未出现边际效应递减的拐点，至少在 AgentNet 数据的观测区间内如此。鉴于密度分析并非本文主线，该结果仅作为补充讨论。

\begin{figure}[htbp]
\centering
\includegraphics[width=0.92\textwidth]{paper01_figures/fig3_验证密度与完成概率.png}
\caption{图 3　验证密度对任务完成概率的预测效应}
\end{figure}


\textbf{图 3} 展示了验证密度与完成概率之间的关系，线性趋势清晰可见，未出现明显的非线性拐点。

\subsubsection{4.6 位置分析：探索性检验}


为检验 H4，本研究对可定位验证关键词的 $N=2{,}512$ 条指令按 25\% / 50\% / 25\% 字符比例划分为开头、中间与末尾三类。三组的完成率分别为 83.87\%、82.01\% 与 81.02\%，卡方检验不显著（χ² = 1.894，$p=.388$，Cramer's V = .028）。成对比例 Z 检验经 Bonferroni 校正后亦均未达到显著水平。

因此，H4 未获支持。在本数据集中，验证表述的"存在"比其"位置"对任务成功更具预测力——这一发现本身具有理论意义，提示终止判准的效能可能更取决于是否提供了验证信息，而非该信息在指令中的位置安排。

\begin{figure}[htbp]
\centering
\includegraphics[width=0.92\textwidth]{paper01_figures/fig4_位置完成率.png}
\caption{图 4　验证表述位于不同位置时的任务完成率}
\end{figure}


\textbf{图 4} 以上述三组为基础，直观展示了不同验证位置下完成率的微小差异及其非显著性。

\subsubsection{4.7 稳健性检验}


\textbf{表 5} 汇总了稳健性检验结果。针对冗余比例，由于约 30\% 任务的冗余比例为 0，本研究采用两部分模型思路：全体样本 OLS、正样本 OLS、是否冗余 Logistic 以及 log(1 + 冗余) OLS。结果显示，`has\_completion\_criteria` 在正样本 OLS 中不显著（β = 0.0007，$p=.740$），在是否冗余 Logistic 中显著为负（β = −0.298，$p<.001$），在其他设定中为弱负向。

上述模式提示一个审慎的结论：H1c（终止阶段冗余减少）未在所有模型设定中获得一致支持；冗余行为可能更受任务领域与模型策略影响，完成标准的关联效应在特定模型规格下才显著。替代密度定义（仅 V1 + V2）与替代位置分段（等分三段）均不改变主结论，进一步验证了 H3、H4 结果的稳健性。


\begin{table}[htbp]
\centering
\caption{表 5：稳健性检验：冗余比例、替代密度与替代位置定义}
\small
\adjustbox{max width=\textwidth}{
\begin{tabular}{l r r l r}
\toprule
模型/设定 & $N$ & 关键统计量 & `has\_completion\_criteria` 效应 & $p$ \\
\midrule
OLS（全体冗余比例） & 20,484 & $R^2=.041$ & β = −0.0057 [−0.0092, −0.0023] & .001 \\
OLS（冗余比例 > 0） & 14,222 & $R^2=.130$ & β = 0.0007 [−0.0032, 0.0045] & .740 \\
Logistic（是否冗余） & 20,484 & 伪 $R^2=.006$ & β = −0.298 [−0.385, −0.212] & <.001 \\
OLS（log(1 + 冗余比例)） & 20,484 & $R^2=.038$ & β = −0.0053 [−0.0084, −0.0023] & <.001 \\
替代密度（V1+V2）二次项 & 20,484 & — & 二次项 $p=.074$ & .074 \\
替代位置（等分三段） & 2,512 & χ² = 3.926 & — & .140 \\
\bottomrule
\end{tabular}
}
\end{table}

\subsection{5 讨论}


\subsubsection{5.1 完成标准作为终止判准：过程层面的证据}


本研究将用户指令中的完成标准/验证条件重新框架为任务型交互智能体的"终止判准"，并基于大规模真实交互轨迹提供了过程层面的证据。主结果显示，含完成标准的任务不仅完成率更高、执行步数更少，而且在终止阶段具有更低的冗余比例、错误比例及异常短轨迹率。这三条证据链——结果（完成率）、总量（步数）与过程（终止阶段行为）——相互印证，共同支持如下解释：当用户指令提供明确的验证信息时，智能体更可能将自然语言目标转译为可操作的停止条件，从而在接近目标时与更少的冗余探索及错误尝试相关联。

需要再次强调，本研究为观察性研究，无法确立因果关系。完成标准的存在可能与任务本身较简单、用户表达能力较强、或模型对特定表述更敏感等未观测变量相关。因此，本文使用"与……相关联""提示""表明"等关联性语言，而非"导致""提升"等因果性表述。

\subsubsection{5.2 理论解释}


上述发现可置于三条理论线索的框架中获得统一解释。

从\textbf{言语行为理论}视角看，完成标准与验证动作使指令的语力更加明确：智能体能够将"达到某状态"或"执行某验证动作"作为行为终止的判据，这在本质上是从指令中提取了更丰富的言外之力信息（Austin, 1962; Searle, 1969）。

\textbf{关联理论}进一步解释为何明确验证信息可减少执行步数：验证关键词直接标记了智能体在规划阶段应建立的子目标——"确认目标状态已达成"，从而使智能体以更低的认知努力即可确认任务完成，避免冗余探索（Sperber \& Wilson, 1995）。值得注意的是，验证密度未呈现倒 U 形效应，这可能意味着在当前 AgentNet 数据的验证密度范围内（绝大多数任务验证密度极低），验证信息的边际效应尚未进入递减区间。

\textbf{行动终止理论}则解释了终止阶段行为差异的机理：完成标准为智能体提供了内部终止规则——"当验证动作确认了目标状态时停止"——减少了对"步数足够多""状态变化趋于平缓"等外部启发式规则的依赖，从而抑制了过冲行为（Gigerenzer et al., 1999; Payne et al., 1993）。异常短轨迹率的降低进一步支持了这一解释：缺乏内部终止规则的任务更容易在智能体尚未充分探索时即触发启发式停止。

\subsubsection{5.3 对智能体设计的实践启示}


本研究的发现对用户指令撰写与智能体终止机制设计具有双向启示。

对\textbf{用户侧}而言，在指令中明确加入"请确认……是否完成""确保……处于……状态"等验证性表述，有助于智能体准确判断终止时机。实践中，可建议用户在描述完操作步骤后，附加一到两句验证性语句，直接指明期望的终态。

对\textbf{智能体设计者}而言，可在规划模块中显式识别验证关键词并将其编码为候选终止条件。具体而言：（1）在指令解析阶段，将验证关键词映射为内部子目标标记；（2）在执行循环的每一步之后，检查当前状态是否满足任一被编码的终止条件；（3）当用户指令缺乏验证信息时，智能体可主动询问用户以建立共同基础（Clark \& Brennan, 1991），例如："请确认：任务完成的标准是什么？"。

此外，由于验证密度未呈现倒 U 形，当前结果不支持对验证信息数量的过度担忧。换言之，在提示工程实践中，鼓励用户加入验证信息至少在 AgentNet 所代表的任务范围内风险较低。

\subsubsection{5.4 局限性与未来方向}


本研究存在以下局限，需在解读结论时加以考虑。

\textbf{第一，研究设计的观测性。} 作为对公开交互轨迹的二次分析，本研究无法操纵完成标准的存在与否，难以直接推断因果。完成标准可能与用户表达能力、任务复杂度等未观测变量相关。未来研究可通过实验操纵验证显性度——例如，对同一任务随机分配含/不含完成标准的指令——以建立因果推断。

\textbf{第二，验证类型的标注精度。} 当前基于关键词规则的标注方法较为粗糙：状态目标式（V2）与条件验证式（V3）的样本量极小（分别为 5 条和 1 条），限制了对 H2 的检验。未来可采用依存句法分析或大语言模型辅助标注，提升验证类型的识别精度和覆盖率。

\textbf{第三，领域与模型的异质性。} 当前分析未区分应用领域（如办公自动化、网页浏览、文件管理）与智能体模型版本，这些因素可能调节完成标准的效应。未来可引入这些变量进行分层分析。

\textbf{第四，终止阶段指标的效度。} 本文的终止阶段指标基于规则化的步骤标注（`last\_step\_redundant`、`last\_step\_correct`），其有效性受原始标注质量与规则定义影响。未来可结合人工复核或更细粒度的行为编码提升指标精度，例如区分"验证性冗余"（有意义的确认步骤）与"探索性冗余"（无目标的试探步骤）。

\subsection{6 结论与展望}


本研究基于 AgentNet 数据集，从言语行为理论、关联理论与行动终止理论出发，将用户指令中的完成标准/验证条件重新概念化为任务型交互智能体的"终止判准"，并实证检验了其与执行效率的关联。主要发现如下：

\begin{enumerate}
\item[1.] 含完成标准的指令与更高的任务完成率、更少的执行步数以及更低的终止阶段冗余/错误率相关联，为"完成标准作为终止判准"的理论框架提供了过程层面的经验证据；
\item[2.] 验证密度与完成率之间未发现稳健的倒 U 形关系，在当前数据中验证信息的增加主要表现为线性正向关联；
\item[3.] 验证表述的"存在性"比其"位置"对任务成功更具预测力，提示智能体终止机制的设计应优先关注验证信息的提供与否，而非其指令内位置布局。
\end{enumerate}


上述发现为理解自然语言指令如何塑造智能体终止行为提供了新的理论与实证基础。在实践层面，建议用户指令明确包含验证性表述，建议智能体系统将验证关键词编码为内部终止条件。

未来研究可从以下方向深化：（1）采用人工标注或依存句法分析提升验证类型识别的精度；（2）通过实验操纵验证显性度以建立因果推断；（3）探讨完成标准与约束密度、任务复杂度三者之间的交互效应——这一方向将在本系列研究的后续论文中展开；（4）将验证信息纳入智能体规划与终止模块，设计可解释的终止决策机制。


\newpage

\chapter{第2部分　约束何时成为负荷：任务型交互智能体用户指令中约束密度的条件性非线性效应——基于 AgentNet 交互轨迹的实证研究}


\subsection{1 引言}


\subsubsection{1.1 从验证到约束：问题域的延伸}


在前一项研究中（见本系列 Paper 01），本文将用户指令中的完成标准/验证条件重新概念化为任务型交互智能体的"终止判准"，并基于 AgentNet 轨迹数据提供了实证证据：含完成标准的指令与更高的任务完成率、更少的执行步数及更优的终止阶段行为相关联。该研究的核心逻辑在于——明确的验证信息为智能体提供了可操作的停止信号，降低了终止不确定性。

然而，用户指令中还存在另一类具有潜在行为约束力的语言成分：约束表达（constraint expressions）。不同于验证信息（告诉智能体"何时算完成"），约束表达告诉智能体"不能做什么""必须做什么"或"只能做什么"——换言之，它通过对可接受动作集合的限定来塑形智能体的行为空间。验证信息界定任务的\textbf{终点}，约束信息界定任务的\textbf{边界}。二者共同构成用户指令中语用信息的两个核心维度，但约束表达的效应逻辑可能更为复杂：验证信息主要作用于终止阶段，而约束信息贯穿规划、执行与验证全过程，其密度效应可能呈现非线性的剂量—反应关系。

\subsubsection{1.2 约束密度的核心问题}


近年以 LLM 为核心的任务型交互智能体在网页浏览、桌面办公与跨应用工作流中展现出强大潜力（Li et al., 2024; Yang et al., 2022; Yuan et al., 2023），用户指令作为人机协作的"语言入口"，其信息结构与精确度直接影响智能体的行为轨迹（Lasecki et al., 2013; Liang et al., 2020）。

在指令设计中，约束表达（如"不要……""必须……""只能……""保持……不变"）是一种常见却未被充分理解的语用手段。从直觉上看，约束通过限定可接受动作集合来提升智能体输出的精确性。然而，一个关键问题悬而未决：\textbf{约束是否越多越好？} 当约束密度超过某一阈值时，是否会因规划、验证与冲突消解成本的上升而反而降低完成率？换言之，约束密度与任务完成率之间是否存在倒 U 形关系？如果存在，这一关系的拐点是否因任务情境而异？

既有文献多将约束视为二元特征（有/无），或将其纳入更宽泛的"指令明确性"框架中讨论（Goedert et al., 2013; Schatz, 2003）。少量研究开始关注约束的语义类型，但对其密度（density）的剂量—反应关系缺乏系统检验。本研究的核心问题是：\textbf{用户指令中的约束密度如何调节任务型交互智能体的动作空间？在什么条件下约束发挥边界作用，又在什么条件下转化为认知负荷？}

\subsubsection{1.3 研究边界与核心论点}


本研究以 AgentNet 公开数据集中的 GUI 轨迹为经验基础（Zhang et al., 2024）。AgentNet 记录了智能体在真实桌面与网页应用环境中执行用户指令的完整操作序列，因此本文所讨论的"任务型交互智能体"特指以 GUI 操作为边界、以自然语言指令为输入的通用型数字智能体。

需要明确的是，AgentNet 数据属于观测性轨迹数据：研究者只能观察到用户"写下了什么指令"以及智能体"最终是否完成"，无法观测用户撰写指令时的真实意图与先验知识。因此，本文所有结论均应在"关联性"而非"因果性"的框架下解读。

本研究的核心论点是：约束在任务型交互智能体用户指令中具有\textbf{双重角色}。一方面，适度约束能够减少目标歧义、压缩搜索空间、降低规划分支数，从而提升完成概率；另一方面，当约束密度过高时，智能体需要额外的计算与推理资源来检查约束一致性、处理约束冲突、在局部可行解之间进行权衡，这会增加认知与计算负荷，最终可能抵消边界带来的收益。因此，约束密度与任务完成率之间的关系更可能呈现一种\textbf{条件性非线性模式}，而非普适的线性或全局倒 U 形规律。具体而言：

\begin{itemize}
\item \textbf{低密度}：约束不足，动作空间过大，目标边界模糊，智能体易产生冗余探索；
\item \textbf{中密度}：约束提供清晰边界，歧义降低，完成率提升；
\item \textbf{高密度}：约束数量超过智能体的规划与验证容量，冲突风险上升，完成率下降。
\end{itemize}


更为关键的是，这一非线性关系的拐点与曲率取决于任务情境。高难度任务、长指令任务以及包含多个子任务的任务，其信息需求更高、规划深度更大，因而更容易因约束过载而受损。

\subsection{2 理论基础与研究假设}


本研究的理论框架由两条相互交织的线索构成：语用学理论解释约束表达为何具有信息价值，认知负荷理论解释该价值为何存在边界。两条线索共同指向一个核心推论：约束密度是条件性资源，其最优值随任务情境而变化。

\subsubsection{2.1 语用学视角：合作原则与关联理论中的约束}


从语用学角度看，用户指令本质上是一种以行动达成为目的的指令式言语行为（directive speech act）（Austin, 1962; Searle, 1969）。Grice（1975）的合作原则（Cooperative Principle）指出，成功交际需遵循量的准则（信息量适中）、质的准则（真实可信）、关系准则（内容相关）与方式准则（表达清晰）。在 GUI 指令中，约束表达是方式准则与量的准则共同作用的产物：它通过明确"不能做什么""必须做什么"来降低歧义，但过量的约束又可能违反量的准则，造成信息冗余。

Sperber 与 Wilson（1995）的关联理论（Relevance Theory）进一步为理解约束密度的边界提供了精细的分析工具。该理论强调，语境效果的获取需要付出认知加工成本；听话者（或信息处理系统）会在认知投入与语境收益之间进行权衡。将这一框架迁移到人机场景可推知：智能体在解码指令时，需将每条约束纳入规划与验证；当约束密度超过一定阈值，新增的语境收益可能无法抵消递增的认知与计算成本。Yeung（2016）在工作指令与培训研究中提出的"过犹不及"假说（too-much-of-a-good-thing）也为这一非线性预期提供了来自组织行为学的独立支持。

值得与 Paper 01 进行理论层面的对比：验证信息（完成标准）主要贡献于终止效率，其"过犹不及"风险相对较低——Paper 01 中验证密度未呈现倒 U 形即为佐证。约束信息则不同：它不仅影响终止判断，更影响全局动作选择空间，因此其"过犹不及"风险更为突出——这正是本研究关注约束密度非线性效应的理论动因。

\subsubsection{2.2 认知负荷与边界条件框架}


认知负荷理论（Cognitive Load Theory）为理解约束密度的条件效应提供了第二视角。Sweller（1988）指出，人类工作记忆容量有限，当外部信息结构超出加工容量时，学习或问题解决绩效会下降。类比到智能体系统，尽管 LLM 在推理阶段并不受限于人类工作记忆，但其规划与验证模块仍受上下文长度、注意力分配与冲突消解能力的约束。高密度约束可能通过以下机制增加系统层面的"负荷"：

\begin{enumerate}
\item[1.] \textbf{规划负荷}：每条约束都需被纳入动作序列生成，约束越多，可行动作组合越复杂；
\item[2.] \textbf{验证负荷}：执行过程中需反复确认当前动作是否违反任一约束；
\item[3.] \textbf{冲突消解负荷}：当多条约束相互冲突时，智能体需进行优先级排序或折中推理；
\item[4.] \textbf{信息检索负荷}：长指令中的高密度约束会稀释关键目标信息，增加语义解析难度。
\end{enumerate}


基于上述理论，本研究将约束密度的效应置于\textbf{边界条件}（boundary conditions）框架下：不存在放之四海而皆准的最优密度，效应方向与强度取决于任务需要多少边界、智能体能承受多少验证负荷。这一框架的核心方法论含义是：寻找全局最优密度可能是一个错误的问题；正确的问题应是——在何种条件下，约束密度呈现倒 U 形效应？

\subsubsection{2.3 研究假设}


综合上述理论，本研究提出以下假设：

\textbf{H1（全局效应弱化）}：在混合所有任务的完整样本中，约束密度与任务完成率之间不存在显著的全局倒 U 形关系。该假设的提出基于数据分布的现实：AgentNet 中约 93\% 的任务不含任何约束关键词，约束密度呈高度零膨胀分布，全局模型会被大量零值稀释，难以检测到统一的曲率。H1 并非否定约束的价值，而是否定"存在一个适用于所有任务的统一最优密度"这一过于简化的命题。

\textbf{H2（条件倒 U 形）}：在任务难度较高、指令长度较长或包含多个子任务的任务子样本中，约束密度与完成率呈显著倒 U 形关系。此时期约束的边界收益与验证成本之间的权衡最为明显——这正是本研究理论框架的核心预测。

\textbf{H3（子类型异质性，探索性）}：不同约束子类型（禁止型、保留型、要求型、限制型）的最优密度可能存在差异。例如，禁止型约束可能因直接排除动作而具有较低的过载阈值。该假设为探索性，因各子类型样本量较小、类型重叠度高。

\textbf{H4（任务难度调节曲率）}：任务难度显著调节约束密度的二次效应，即二次项与任务难度的交互项为负且显著。难度越高，约束密度的倒 U 形曲率越强——这一假设是对 H2 的精细化，旨在检验调节效应是否具体作用于曲率而非线性斜率。

\subsection{3 研究方法}


\subsubsection{3.1 数据来源与样本构成}


本研究使用 AgentNet 公开数据集的预处理版本（Zhang et al., 2024）。原始样本包含 22,811 条 GUI 任务轨迹。经完整案例分析（listwise deletion）后，主模型纳入 20,484 条任务。其中，指令文本中包含至少一个约束关键词的原始任务为 1,465 条（占 6.42\%），进入回归模型的约束任务为 1,261 条。因变量为二值变量 `task\_completed`（1 = 完成，0 = 未完成），由轨迹层面的执行结果决定。

本研究的样本构成具有一个不可忽视的结构特征——约束密度的高度零膨胀分布（约 93\% 的任务零约束）——这意味着整体样本中约束密度的变异主要由"有/无约束"的二元差异驱动，而非由密度的连续变化驱动。这一特征直接决定了分析策略的选择：必须将"约束有无效应"与"约束密度效应"分开建模。

\subsubsection{3.2 约束词分类与密度操作化}


参考语用学中情态、言外之力与限定表达研究，本研究将约束表达划分为五类（见表 1）：禁止型（C1）、保留型（C2）、要求型（C3）、限制型（C4）以及混合型（C5，同时包含两类及以上约束）。关键词匹配采用 Python 正则表达式，按不区分大小写、词边界方式匹配；短语优先于单词计数以避免重复。

约束密度（`constraint\_density`）定义为约束词绝对频次除以指令词数再乘以 100，以百分比表示；约束丰富度（`constraint\_richness`）定义为同时出现的不同约束类型数（0–4）。


\begin{table}[htbp]
\centering
\caption{表 1：约束表达的类型学定义与关键词}
\small
\adjustbox{max width=\textwidth}{
\begin{tabular}{l l l}
\toprule
类型 & 定义 & 关键词 \\
\midrule
C1 禁止型 & 禁止某些操作或行为 & don't, do not, never, avoid, without, cannot, should not, must not \\
C2 保留型 & 要求保持/保留某状态 & keep, maintain, stay, preserve, leave \\
C3 要求型 & 要求必须完成某事 & must, should, need to, have to, required \\
C4 限制型 & 限定范围、对象或方式 & only, just, merely, exclusively \\
C5 混合型 & 同时包含两类及以上约束 & — \\
\bottomrule
\end{tabular}
}
\end{table}

\subsubsection{3.3 三阶段分析框架}


为系统回答"约束何时成为负荷"，本研究采用逐层递进的三阶段分析框架。

\textbf{Stage 1：约束有无效应。} 在全样本（$N=20{,}484$）中估计 Logistic 模型：

\[
\text{logit}(P(\text{完成})) = \beta_0 + \beta_1 \cdot \text{has\_constraint\_words} + \gamma \cdot \text{task\_difficulty} + \delta \cdot \text{word\_count} + \eta \cdot \text{os}
\]


该阶段回答一个基础性问题：约束表达的"存在"本身是否具有预测价值？结果以比值比（OR）及 95\% Wald 置信区间报告。

\textbf{Stage 2：约束任务内的全局密度效应。} 将样本限制为包含至少一个约束词的任务（$N=1{,}261$），估计二次 Logit 模型：

\[
\text{logit}(P(\text{完成})) = \beta_0 + \beta_1 \cdot \text{density} + \beta_2 \cdot \text{density}^2 + \gamma \cdot \text{task\_difficulty} + \delta \cdot \text{word\_count} + \eta \cdot \text{os}
\]


若 $\beta_2<0$ 且显著，则支持全局倒 U 形假说。拐点由 $-\beta_1/(2\beta_2)$ 计算，使用 Delta 法估计 95\% CI。

\textbf{Stage 3：条件模型。} 通过交互项与分样本模型考察任务难度、指令长度、子任务数量等情境变量的调节作用。交互模型形式为：

\[
\text{logit}(P(\text{完成})) = \beta_0 + \beta_1 d + \beta_2 d^2 + \beta_3 M + \beta_4 dM + \beta_5 d^2 M + \text{controls}
\]


其中 $d$ 为约束密度，$M$ 为调节变量。分样本模型按调节变量的中位数或二分取值分别估计二次 Logit。这一阶段的理论目标是：识别约束密度效应的边界条件，即确定"约束何时成为负荷"的具体情境。

\subsubsection{3.4 辅助检验与数据特征说明}


本研究采用 Dawson（2014）U 形检验判断观测区间两端斜率是否符合倒 U 形预期，使用 Cohen's $f^2$ 评估二次项的解释增量，并以绝对频次、丰富度、限制性立方样条（RCS）及操作系统分样本进行稳健性检验。模型比较以 AIC/BIC 与 McFadden 伪 $R^2$ 为准则。

关于测量有效性，需说明以下几点。其一，约束识别基于关键词规则，具有可解释性与可重复性，但无法识别隐含约束或复杂否定。其二，同一关键词在不同语境中可能承载不同语用功能（例如 "just" 既可是限制词，也可是语气副词），规则匹配会产生一定误判。其三，约束密度以词数为分母，极短指令中单一约束可能产生极端密度值（如 20\%），好在高密度观测的稀缺性（95\% 分位数仅为 1.74\%）使极端值对整体推断的影响有限。尽管存在上述局限，规则检测在现有大规模标注数据稀缺的条件下仍是一种可行且透明的操作化方法。

\subsection{4 结果}


\subsubsection{4.1 约束类型与密度分布}


表 2 展示了约束类型的分布特征。一个显著模式是约束表达的极度稀缺性：无约束任务占 93.58\%，即超过九成的用户指令未包含任何显式约束关键词。在含约束的少数任务中，限制型约束（C4）最为常见（2.01\%），其次为禁止型（1.53\%）、保留型（1.31\%）与要求型（1.10\%），混合型（C5）不足 0.5\%。


\begin{table}[htbp]
\centering
\caption{表 2：约束类型分布（原始样本 $N=22{,}811$）}
\small
\adjustbox{max width=\textwidth}{
\begin{tabular}{l r r}
\toprule
约束类型 & 任务数 & 占全部样本比例（\%） \\
\midrule
无约束 & 21,346 & 93.58 \\
C1 禁止型 & 350 & 1.53 \\
C2 保留型 & 298 & 1.31 \\
C3 要求型 & 251 & 1.10 \\
C4 限制型 & 458 & 2.01 \\
C5 混合型 & 108 & 0.47 \\
\bottomrule
\end{tabular}
}
\end{table}

\begin{figure}[htbp]
\centering
\includegraphics[width=0.92\textwidth]{paper02_figures/fig1_约束类型分布.png}
\caption{图 1　用户指令中约束类型的分布}
\end{figure}


\textbf{图 1} 直观呈现了约束类型的极度不均衡分布——无约束任务构成绝对主体，各约束子类型均属低频类别。这一分布特征对于后续分析的统计功效具有重要影响：子类型层面的推断将受到样本量的严格限制。

表 3 报告了约束密度的描述性统计。整体密度呈高度右偏与零膨胀：均值仅 0.21\%，中位数为 0，95\% 分位数为 1.74\%，最大值达 20.0\%。这意味着绝大多数用户指令在措辞上极为简洁，显式约束表达并非用户自然书写的默认特征。


\begin{table}[htbp]
\centering
\caption{表 3：约束密度的描述性统计（原始样本 $N=22{,}811$）}
\small
\adjustbox{max width=\textwidth}{
\begin{tabular}{l r}
\toprule
统计量 & 约束密度（\%） \\
\midrule
样本量 & 22,811 \\
均值 & 0.206 \\
标准差 & 0.943 \\
最小值 & 0.000 \\
5\% 分位 & 0.000 \\
25\% 分位 & 0.000 \\
中位数 & 0.000 \\
75\% 分位 & 0.000 \\
95\% 分位 & 1.739 \\
最大值 & 20.000 \\
\bottomrule
\end{tabular}
}
\end{table}

\begin{figure}[htbp]
\centering
\includegraphics[width=0.92\textwidth]{paper02_figures/fig2_约束密度分布直方图.png}
\caption{图 2　约束密度的分布直方图}
\end{figure}


\textbf{图 2} 以双面板呈现了约束密度的分布：左面板为全部任务的密度直方图，清晰展示了零膨胀特征；右面板为仅含约束任务的密度直方图，显示即使在约束子集中，密度亦高度集中于低值区间。

\subsubsection{4.2 Stage 1：约束有无效应}


表 4 上半部分报告了 Stage 1 模型结果。在控制任务难度、指令长度与操作系统后，`has\_constraint\_words` 的回归系数为 0.152，对应 OR = 1.165（95\% CI [1.075, 1.262]，$z=3.713$，$p<0.001$）。这一结果表明，在完整样本中，包含至少一个约束词的指令与更高的任务完成率存在显著正相关——平均而言，含约束词指令的完成几率比不含约束词的指令高约 16.5\%。

该结果确立了一个重要的基线：约束表达的"存在"本身具有预测价值。然而，它并未回答本研究关心的核心问题——约束的密度效应。二元效应的显著性并不必然意味着密度的连续增加会持续改善绩效；事实上，这正是 Stage 2 所要检验的关键区分。

\subsubsection{4.3 Stage 2：约束任务内的全局密度效应}


将样本限制为包含至少一个约束词的任务（$N=1{,}261$）后，估计二次 Logit 模型。表 4 下半部分报告了结果：约束密度的线性项显著为负（$b=-0.267$，$p=0.034$），但二次项为正且未达传统显著性水平（$b_2=0.016$，$p=0.083$）。拐点估计值为 8.46\%（95\% CI [4.58, 12.35]），但鉴于二次项不显著，该拐点的统计证据不足。Dawson U 形检验显示，在密度下限（0.57\%）处斜率为 −0.249（$p=0.032$），上限（20.0\%）处斜率为 0.364（$p=0.150$），未同时满足倒 U 形的两端斜率条件。


\begin{table}[htbp]
\centering
\caption{表 4：Stage 1 与 Stage 2 Logistic 回归结果}
\small
\adjustbox{max width=\textwidth}{
\begin{tabular}{l l r l r}
\toprule
变量 & Stage 1：完整样本（$N=20{,}484$） &  & Stage 2：约束任务子样本（$N=1{,}261$） &  \\
\midrule
 & Coef (SE) & OR [95\% CI] & Coef (SE) & OR [95\% CI] \\
Intercept & 1.851 (0.086)*** & 6.364 [5.381, 7.528] & 3.163 (0.650)*** & 23.645 [6.613, 84.546] \\
os [win\_mac] & −0.589 (0.060)*** & 0.555 [0.493, 0.624] & −1.057 (0.254)*** & 0.347 [0.211, 0.572] \\
has\_constraint\_words & 0.152 (0.041)*** & 1.165 [1.075, 1.262] & — & — \\
constraint\_density & — & — & −0.267 (0.126)* & 0.765 [0.598, 0.980] \\
I(constraint\_density²) & — & — & 0.016 (0.009) & 1.016 [0.998, 1.034] \\
task\_difficulty & −0.016 (0.013) & 0.984 [0.958, 1.010] & −0.015 (0.064) & 0.985 [0.869, 1.116] \\
word\_count & 0.002 (0.001) & 1.002 [1.000, 1.004] & −0.004 (0.005) & 0.996 [0.987, 1.006] \\
拐点 & — & — & 8.46 [4.58, 12.35] & — \\
\bottomrule
\end{tabular}
}
\end{table}

\emph{注：} $^*p<0.05$，$^{**}p<0.01$，$^{***}p<0.001$。Stage 2 密度下限为 0.57\%，上限为 20.0\%。

这一结果支持 H1：在约束任务内部，不存在稳健的全局倒 U 形密度效应。换言之，一旦用户决定在指令中加入约束，单纯增加约束密度并未显著改善或显著降低完成率；密度效应在约束任务整体上趋于平缓。

\begin{figure}[htbp]
\centering
\includegraphics[width=0.92\textwidth]{paper02_figures/fig3_约束密度与完成率关系.png}
\caption{图 3　约束密度与任务完成率的分箱观测值及二次拟合}
\end{figure}


\textbf{图 3} 以分箱观测值与二次拟合曲线展示了约束密度与完成率的关系。二次拟合曲线在大部分密度区间内近乎平坦，仅在极高密度处略微上行——这与 H1 的预期一致，即全局曲率不显著。

Stage 2 的非显著结果本身构成一个有意义的方法论信号：它提示我们，约束密度的效应可能不是全局性的，而是高度情境依赖的。这一信号自然引导分析进入 Stage 3——条件效应检验。

\subsubsection{4.4 Stage 3：条件模型——任务难度、指令长度与子任务数量}


\paragraph{4.4.1 任务难度的调节作用}


表 5 报告了交互效应模型与分样本二次模型结果。在交互模型中，约束密度二次项与任务难度的交互项显著为负（$b=-0.0091$，$p=0.019$），而密度线性项与难度的交互项不显著（$b=0.036$，$p=0.180$）。这意味着任务难度调节的是约束密度的\textbf{曲率}而非其线性斜率——随着难度增加，密度—完成率关系从近乎平坦逐渐弯曲为倒 U 形。

按任务难度中位数（6）分样本后，模式更为清晰：高难度任务子样本（$N=8{,}413$）的二次项显著为负（$b_2=-0.039$，$p=0.019$），拐点约 1.93\%（95\% CI [1.03, 2.82]）；低难度任务子样本（$N=12{,}071$）的二次项不显著（$b_2=0.007$，$p=0.218$）。该结果支持 H2 与 H4，为本研究的核心理论预测——约束密度效应取决于任务情境——提供了关键的实证证据。

\paragraph{4.4.2 指令长度的调节作用}


按词数中位数（30）分样本后，长指令任务子样本（$N=10{,}211$）呈现显著倒 U 形（$b_2=-0.041$，$p=0.025$），拐点约 1.69\%（95\% CI [0.85, 2.52]）；短指令任务子样本（$N=10{,}273$）的二次项不显著（$b_2=0.008$，$p=0.212$）。长指令与高密度约束叠加时，信息稀释与验证负荷效应更加明显——这与关联理论的预期一致：在信息总量已经较大的情况下，新增约束的边际语境收益递减，而加工成本递增。

\paragraph{4.4.3 子任务数量的探索性条件效应}


除任务难度与指令长度外，本研究进一步考察了 `has\_multiple\_subtasks` 的调节作用。在包含多个子任务的任务子样本（$N=13{,}109$）中，约束密度呈现显著倒 U 形（$b_2=-0.045$，$p=0.006$），拐点约 1.87\%（95\% CI [1.23, 2.51]）；单个子任务子样本（$N=7{,}375$）中二次项不显著（$b_2=0.010$，$p=0.159$）。这一发现与任务难度、长指令结果形成高度一致的三角验证：三种操作化方式均指向同一结论——任务结构越复杂，约束密度的倒 U 形效应越显著。


\begin{table}[htbp]
\centering
\caption{表 5：约束密度的条件效应模型}
\small
\adjustbox{max width=\textwidth}{
\begin{tabular}{l l l r r}
\toprule
模型 & 变量 & Coef (SE) & OR [95\% CI] & $p$ \\
\midrule
交互模型（$N=20{,}484$） & constraint\_density & −0.167 (0.157) & 0.846 [0.622, 1.152] & 0.288 \\
 & task\_difficulty & −0.014 (0.014) & 0.986 [0.960, 1.012] & 0.287 \\
 & constraint\_density × task\_difficulty & 0.036 (0.027) & 1.037 [0.984, 1.093] & 0.180 \\
 & I(constraint\_density²) & 0.036 (0.020) & 1.037 [0.998, 1.078] & 0.064 \\
 & \textbf{I(constraint\_density²) × task\_difficulty} & \textbf{−0.009 (0.004)} & \textbf{0.991 [0.983, 0.998]} & \textbf{0.019} \\
 & word\_count & 0.003 (0.001)* & 1.003 [1.001, 1.005] & 0.011 \\
分样本 & 分组 & $N$ & $b_2$ & $p$ \\
 & 低难度 & 12,071 & 0.007 & 0.218 \\
 & \textbf{高难度} & \textbf{8,413} & \textbf{−0.039} & \textbf{0.019} \\
 & 短指令 & 10,273 & 0.008 & 0.212 \\
 & \textbf{长指令} & \textbf{10,211} & \textbf{−0.041} & \textbf{0.025} \\
 & 单个子任务 & 7,375 & 0.010 & 0.159 \\
 & \textbf{多个子任务} & \textbf{13,109} & \textbf{−0.045} & \textbf{0.006} \\
\bottomrule
\end{tabular}
}
\end{table}

\begin{figure}[htbp]
\centering
\includegraphics[width=0.92\textwidth]{paper02_figures/fig4_revised_条件二次项系数森林图.png}
\caption{图 4　条件模型中约束密度二次项系数的森林图}
\end{figure}


\textbf{图 4} 以森林图形式呈现了不同分样本中约束密度二次项系数及其 95\% CI。高难度、长指令与多子任务三个分样本的二次项系数均显著为负且落在相似区间，而对应的低水平分样本均不显著——这一"条件性显著"模式是 H2 最直观的可视化证据。

\begin{figure}[htbp]
\centering
\includegraphics[width=0.92\textwidth]{paper02_figures/fig5_难度调节效应.png}
\caption{图 5　任务难度对约束密度效应的调节}
\end{figure}


\textbf{图 5} 聚焦于任务难度的调节效应，以多面板形式展示了低、中、高三种难度水平下约束密度与完成率的关系。高难度任务中约束密度的倒 U 形最为清晰，最优密度区间约 1.7\%–1.9\%；低难度任务中密度效应近乎平坦。该图与图 4 的森林图互为补充——图 4 从系数估计角度呈现条件效应的统计证据，图 5 从预测概率角度呈现条件效应的实际形态。

\paragraph{4.4.4 Johnson-Neyman 风格的调节区间分析}


为更精确刻画任务难度的调节区间，本研究基于交互模型计算了不同任务难度水平下约束密度二次项的曲率及其标准误。附录表 9 报告了完整的简单斜率与曲率序列。结果显示，当任务难度约 ≥ 6 时，曲率显著为负（$p<0.05$）；难度低于该阈值时，曲率不显著或接近零。

\begin{figure}[htbp]
\centering
\includegraphics[width=0.92\textwidth]{paper02_figures/fig_johnson_neyman_curvature.png}
\caption{图 6　任务难度对约束密度曲率的 Johnson-Neyman 风格调节}
\end{figure}


\textbf{图 6} 以可视化方式呈现了 Johnson-Neyman 风格的边界：横轴为任务难度（0–10），纵轴为约束密度的曲率估计值。红色阴影区域标示曲率显著为负的难度区间（约 ≥ 6），蓝色区域标示曲率不显著的区间。该图直观揭示了"约束何时成为负荷"的条件边界。

\begin{figure}[htbp]
\centering
\includegraphics[width=0.92\textwidth]{paper02_figures/fig_conditional_predicted_probs.png}
\caption{图 7　不同任务难度水平下约束密度的预测完成概率}
\end{figure}


\textbf{图 7} 进一步展示了低（3 分）、中（6 分）、高（9 分）三种难度水平下约束密度的预测完成概率曲线。高难度任务的曲线呈现清晰的倒 U 形——完成概率在密度约 1.9\% 处达到峰值，随后下行；中难度任务出现微弱的弯曲趋势；低难度任务的曲线近乎平坦。

\subsubsection{4.5 子类型异质性（探索性）}


表 6 报告了按子类型分别拟合二次模型的结果。四类子类型的二次项系数均未达到显著水平：禁止型 $b_2=-0.0002$（$p=0.986$），保留型 $b_2=0.011$（$p=0.477$），要求型 $b_2=0.014$（$p=0.250$），限制型 $b_2=0.015$（$p=0.356$）。H3 未获支持。


\begin{table}[htbp]
\centering
\caption{表 6：各约束子类型的二次模型结果}
\small
\adjustbox{max width=\textwidth}{
\begin{tabular}{l r l r r r l}
\toprule
子类型 & 非零样本量 & $b_1$ & $p(b_1)$ & $b_2$ & $p(b_2)$ & 拐点 [95\% CI] \\
\midrule
C1 禁止型 & 355 & −0.035 & 0.615 & −0.000 & 0.986 & −102.32 [−11491.34, 11286.70] \\
C2 保留型 & 316 & −0.081 & 0.379 & 0.011 & 0.477 & 3.72 [−0.62, 8.05] \\
C3 要求型 & 224 & −0.182 & 0.029 & 0.014 & 0.250 & 6.74 [−0.47, 13.95] \\
C4 限制型 & 459 & −0.048 & 0.579 & 0.015 & 0.356 & 1.58 [−1.40, 4.56] \\
\bottomrule
\end{tabular}
}
\end{table}

\begin{figure}[htbp]
\centering
\includegraphics[width=0.92\textwidth]{paper02_figures/fig4_子类型拐点对比.png}
\caption{图 8　各约束子类型二次项系数的森林图}
\end{figure}


\textbf{图 8} 以森林图呈现了各子类型二次项系数的估计值与置信区间，所有区间均跨过零线，表明无一子类型达到显著水平。

子类型效应不显著可能源于两方面原因。其一，统计功效不足——各子类型的非零案例仅 224–459 条，且混合型任务的存在进一步稀释了单一类型的独立效应。其二，规则式关键词匹配无法捕捉同一类型内部的语义强度差异（例如 "must" 与 "should" 的强制程度不同）。因此，子类型差异应被视为探索性结果，未来可在更大规模、更高质量标注的数据上重新检验。

\subsubsection{4.6 稳健性检验}


表 7 汇总了多种稳健性检验结果。无论是以绝对频次、丰富度还是 RCS 样条替代二次密度，AIC/BIC 与伪 $R^2$ 均与主模型几乎一致（AIC 范围 20822–20826，伪 $R^2$ 范围 0.0058–0.0059），说明 Stage 2 的全局效应不显著并非操作化方式所致。按操作系统分样本后，Windows/Mac 与 Ubuntu 子样本的二次项均不显著。值得注意的是，RCS 3 节点模型在 200 次迭代后未收敛，提示在大规模零膨胀数据上使用复杂样条需谨慎。


\begin{table}[htbp]
\centering
\caption{表 7：稳健性检验结果}
\small
\adjustbox{max width=\textwidth}{
\begin{tabular}{l r r r r}
\toprule
模型规格 & 样本量 & AIC & BIC & McFadden $R^2$ \\
\midrule
线性密度 & 20,484 & 20,822.04 & 20,861.68 & 0.0058 \\
二次密度 & 20,484 & 20,823.92 & 20,871.49 & 0.0058 \\
绝对频次二次 & 20,484 & 20,824.14 & 20,871.70 & 0.0058 \\
丰富度二次 & 20,484 & 20,822.62 & 20,870.19 & 0.0059 \\
操作系统：Win/Mac & 17,811 & 18,725.74 & 18,764.68 & 0.0002 \\
操作系统：Ubuntu & 2,673 & 1,968.22 & 1,997.67 & 0.0682 \\
RCS（3 节点） & 20,484 & 20,825.99 & 20,881.48 & 0.0058 \\
\bottomrule
\end{tabular}
}
\end{table}

\begin{figure}[htbp]
\centering
\includegraphics[width=0.92\textwidth]{paper02_figures/fig6_模型比较.png}
\caption{图 9　不同模型规格的 AIC 与 BIC 比较}
\end{figure}


\textbf{图 9} 以条形图形式呈现了不同模型规格的 AIC/BIC 比较，所有模型的拟合优度高度接近，进一步确认了全局二次效应不显著的稳健性。

\subsubsection{4.7 探索性过程结果：冗余步数比的非线性模式}


作为机制探索，本研究进一步以总步数（负二项回归）与冗余步数比（OLS）为结果变量，检验约束密度的过程效应。总步数的二次项不显著（$b_2=0.00010$，$p=0.915$），但冗余步数比呈现显著倒 U 形（$b_2=-0.00038$，$p=0.012$）——中等约束密度时冗余探索最多。

这一模式提示约束密度可能通过"探索—验证"权衡影响执行效率：低密度时智能体缺乏方向而探索不足（或随机探索）；中密度时约束引发了更多定向但尚未完成的探索；高密度时约束过于严格，导致探索被压缩。该发现为约束密度的非线性效应提供了过程层面的机制线索，但因属于观测性探索分析，不宜做过强的因果推断。

\subsection{5 讨论}


\subsubsection{5.1 从"全局定律"到"边界条件"：理论贡献的重新定位}


本研究最重要的理论贡献在于将"约束密度是否最优"这一问题从"寻找普适定律"转向"识别边界条件"。整体样本与约束任务子样本均未发现显著的全局倒 U 形效应，这与"过犹不及"假说（Yeung, 2016）的直接推论不一致——如果约束只是单纯地"越多越糟"，那么全局二次项应当显著为负。然而，调节效应分析揭示了一个关键边界：只有在高难度任务、长指令任务或多子任务任务中，约束密度才表现出显著倒 U 形。

这一发现将本研究与 Paper 01 的发现置于互补的理论位置。Paper 01 表明，验证信息（完成标准）与执行效率呈正向关联，且验证密度未呈现倒 U 形——因为验证信息主要作用于终止阶段，认知成本相对有限。Paper 02 则表明，约束信息的效应更加微妙：其价值高度依赖于任务情境；在简单任务中，约束增减的影响微乎其微；在复杂任务中，约束密度存在一个狭窄的"甜蜜点"（约 1.7\%–1.9\%）。两条线索共同指向一个更一般的理论命题：\textbf{用户指令中不同语用成分的效能边界取决于该成分所触发的认知加工类型与加工负荷。}

\subsubsection{5.2 低、中、高约束密度的作用机制}


综合本研究的各项发现，可对约束密度的作用机制进行如下分层描述：

\begin{itemize}
\item \textbf{低密度（< 1\%）}：约束提供的信息边界不足，智能体需要在较大动作空间中搜索，容易产生偏离目标的探索。Stage 1 结果显示，完全不含约束词的指令完成率显著较低，印证了边界缺失的代价。
\item \textbf{中密度（约 1.7\%–1.9\%）}：对于复杂任务，该区间是约束收益最大化的"甜蜜点"。此时约束数量足够压缩动作空间、澄清目标，但又未多到引发严重冲突与验证负荷。值得注意的是，高难度、长指令与多子任务任务的拐点均落在 1.7\%–1.9\% 附近——这一跨操作化的一致性为该区间的实际参考价值提供了交叉验证。
\item \textbf{高密度（> 2\%–3\%）}：约束数量超出智能体的有效处理能力。规划模块需在大量约束下搜索可行路径，验证模块需反复检查每一步是否违规，冲突消解模块需处理相互矛盾的约束——各项成本叠加，最终表现为完成率下降。
\end{itemize}


\subsubsection{5.3 与认知负荷理论和关联理论的对话}


从认知负荷视角看（Sweller, 1988），简单任务本身信息需求较低，额外约束带来的收益有限但成本亦低，因此密度效应不显著；复杂任务需要更多指引，但过量约束可能迅速超出智能体的规划容量，与性能下降相伴随（Schatz, 2003; Goedert et al., 2013）。

关联理论（Sperber \& Wilson, 1995）则进一步解释为何长指令任务对高密度约束更敏感：在信息总量已经较大的情况下，新增约束的边际语境收益递减，而加工成本递增。这一"边际收益递减"逻辑与经济学中的"最优边界"概念形成了有趣的理论共鸣。

\subsubsection{5.4 三阶段框架的方法论意义}


本研究采用的三阶段分析框架——约束有无效应 → 全局密度效应 → 条件效应——为研究大规模观测轨迹数据中的非线性效应提供了一种可推广的方法论模板。该框架的核心设计原则是：先确认基础效应的存在，再检验高阶效应的普遍性，最后通过情境变量识别边界条件。这一递进逻辑有效避免了直接将全局非线性模型应用于零膨胀数据所导致的效应稀释与错误推断。

\subsubsection{5.5 实践建议}


基于研究发现，对任务型交互智能体的提示工程提出以下建议：

\begin{enumerate}
\item[1.] \textbf{高难度任务}：将约束密度控制在 1.7\%–1.9\% 左右。避免在已复杂的任务中堆砌过多否定词、限定词或强制词。
\item[2.] \textbf{长指令任务}：鉴于信息稀释效应，长指令中的约束应更精练，优先保留与核心目标直接相关的约束。
\item[3.] \textbf{多子任务任务}：每个子任务可单独给出关键约束，但全局层面的约束密度仍应控制在 1.7\%–1.9\% 附近。
\item[4.] \textbf{简单任务}：约束密度的容忍区间较宽，但仍需避免极端高密度（如 > 5\%）。
\item[5.] \textbf{接口设计}：可在用户输入指令后提供约束密度反馈，对高密度提示给出精简建议，尤其针对高难度或长任务。
\end{enumerate}


\subsubsection{5.6 局限与未来方向}


本研究存在以下局限。第一，约束检测的规则性——基于关键词的方法无法识别隐含约束与语义强度差异，可能引入测量误差。第二，观测性设计——无法排除任务领域、应用类型、用户能力等未观测混淆；所有效应均为关联性。第三，高密度观测稀缺——95\% 分位数仅为 1.74\%，限制了对高密度区间曲率的统计功效。第四，RCS 3 节点模型未收敛，提示在零膨胀稀疏数据上使用复杂非参数方法需谨慎。

未来研究可从以下方向拓展：（1）使用大语言模型进行更精细的约束语义抽取与人工校验；（2）开展控制实验，随机分配不同约束密度以验证因果关系；（3）构建中介模型，检验冗余步数、验证步数是否在密度与完成率之间发挥中介作用；（4）在多模态智能体与机器人任务中验证结论的跨领域适用性；（5）与本系列 Paper 03 中的语用结构类型学进行交叉分析，探讨约束密度如何与整体语用结构类型交互影响执行效果。

\subsection{6 结论与展望}


本研究基于 AgentNet 大规模 GUI 轨迹数据，系统考察了用户指令中约束密度对任务完成率的非线性影响。主要结论如下：

\begin{enumerate}
\item[1.] 在完整样本中，包含约束词与任务完成率存在显著正相关（OR ≈ 1.165），表明约束表达整体上具有预测价值。
\item[2.] 在约束任务子样本中，不存在稳健的全局倒 U 形密度效应；约束密度的二次项不显著。
\item[3.] 任务难度、指令长度与子任务数量显著调节约束密度的效应：在高难度、长指令或多子任务情境下，约束密度呈现显著倒 U 形，最优密度集中在 1.7\%–1.9\% 附近。
\item[4.] 冗余步数比在中等密度时达到峰值，暗示约束密度可能通过探索—验证权衡影响执行效率。
\end{enumerate}


本研究的理论意义在于将"约束密度最优"问题重新界定为"约束何时成为负荷"的边界条件问题——这一重新框架化将问题从对"普适最优值"的无效寻找中解放出来，转向对"情境条件"的精准识别。实践意义在于为提示工程提供了可量化的密度参考区间。未来研究可在更大规模、更高质量标注数据的基础上，进一步探索约束的语义角色、因果机制与跨任务泛化能力，并将本研究的发现与语用结构类型学（Paper 03）进行整合分析。


\newpage

\chapter{第3部分　任务型交互智能体用户指令的语用结构类型及其执行效应——基于 AgentNet 交互轨迹的潜类别分析}


\subsection{1 引言}


\subsubsection{1.1 从"成分"到"结构"：分析单元的提升}


本系列的前两项研究分别考察了用户指令中两类关键语用成分的独立效应：Paper 01 将完成标准/验证条件重新定位为任务型交互智能体的"终止判准"，提供了其与执行效率正向关联的经验证据；Paper 02 将约束密度界定为"条件性边界资源"，揭示了其与任务完成率之间的非线性关系如何随任务情境而变化。两项研究在各自的分析层面提供了有价值的发现，但均聚焦于\textbf{单一语用成分}——验证信息或约束信息——与执行效果之间的关联。

然而，真实的用户指令并非由孤立语用成分拼凑而成。一条指令可能同时包含操作对象描述、动作顺序说明、文件路径引用、约束条件与完成标准，也可能仅由一个简短的动作短语构成。用户如何将这些信息维度\textbf{组合}为一个整体的语用结构——而非各维度的简单加总——才是理解指令质量的关键。换言之，前两项研究回答了"某一成分是否重要"，本研究则回答一个更为根本的问题：\textbf{用户如何通过自然语言将任务组织为可执行的语用结构？这些结构类型在执行效果上是否存在系统性差异？}

\subsubsection{1.2 指令质量的"组合性"问题}


任务型交互智能体以自然语言指令为"语言入口"，其典型执行链路可概括为：用户任务表达 → 任务表征 → 动作规划 → 环境操作 → 反馈修正 → 结果验证 → 执行终止。在这一链路中，指令不仅是语义输入的载体，更是智能体构建任务表征、设定动作边界与判断完成状态的语用框架（Lasecki et al., 2013; Liang et al., 2020）。

现有研究多从任务完成率、步骤长度、错误率等结果指标评估智能体性能，对输入端指令的语用结构关注不足（Zhang et al., 2024; Li et al., 2024）。提示工程的实践表明，指令的明确性、约束性与结构性影响模型对用户意图的解析精度（Yang et al., 2022; Pal et al., 2022）。然而，一个关键的方法论问题尚未被充分探讨：指令质量的"组合性"——即不同语用维度的\textbf{组织方式}本身是否构成一个独立于各维度简单加总的信息源？

这一问题具有直接的方法论含义。若语用维度的组合方式确实与执行效果存在独立关联，那么仅用显式性评分（将各维度相加）或单一维度（如约束密度）来评估指令质量，在方法论上就是不充分的。这也正是 Paper 01 与 Paper 02 各自所不能回答的问题：验证密度与约束密度是否只是同一潜在语用结构的两个表层指标？不同类型用户指令的语用结构是否存在可识别的"类型"，而非连续谱系上的任意一点？

\subsubsection{1.3 研究问题与定位}


本研究以 AgentNet GUI 交互轨迹为经验基础（Zhang et al., 2024），借鉴语用类型学、言语行为分类学（Austin, 1962; Searle, 1976）与共同基础理论（Clark \& Brennan, 1991），构建包含 11 个二元语用特征的分析框架，并采用潜类别分析（LCA）识别用户指令的潜在语用结构类型。本研究的核心问题有三个递进层次：

\begin{enumerate}
\item[1.] \textbf{存在性问题}：AgentNet 用户指令是否存在可识别的潜在语用结构类型？
\item[2.] \textbf{差异性问量}：不同语用结构类型在执行效果（完成率、过程指标、路径稳定性）上是否存在系统性差异？
\item[3.] \textbf{独立性问题}：语用结构类型的效应是否独立于线性显式性评分？是否受任务难度调节？
\end{enumerate}


研究方法论上的定位是：将分析单元从 Paper 01/02 的"单一语用成分"提升至"整体语用结构"，从而为指令质量的组合性评估提供经验基础。

\subsection{2 理论基础与研究假设}


本研究的理论框架由三条线索交织而成：语用类型学与言语行为分类学界定"有哪些语用维度"；共同基础理论解释"为什么语用结构差异会影响执行"；潜类别分析提供"如何识别类型"的方法论基础。

\subsubsection{2.1 语用类型学与言语行为分类学}


语用类型学关注语言使用者在特定交际场景中对语言形式的选择及其功能后果。Austin（1962）的言语行为理论将语言视为行为，区分言内行为（locutionary act）、言外行为（illocutionary act）与言后行为（perlocutionary act）。Searle（1976）进一步将言外行为划分为断言类、指令类、承诺类、表达类与宣告类五类。在任务型交互智能体场景中，用户指令主要属于指令类言语行为——其核心语力（illocutionary force）在于促使智能体执行特定操作。

然而，指令类言语行为内部存在丰富的结构性变异。有些指令强调动作本身（"点击按钮""输入文字"），有些强调操作对象（"打开那个文件""修改第三张表格"），有些强调约束条件（"只能使用终端""不要修改原文件"），还有些强调完成标准（"确认重命名成功""检查是否在打折"）。这些差异并非程度上的"多一点或少一点"，而是\textbf{组织方式}上的质的不同——它们反映了用户如何将任务映射为语言表达的不同语用策略。

本研究从指令文本中提取的 11 个二元语用特征（是否明确应用、操作对象、目标动作、目标状态、文件引用、对象属性、操作顺序、多子任务、约束词、完成标准、初始状态）覆盖了"谁/什么/怎么做/做到什么程度"四个语用维度，构成一个系统性的语用特征空间。不同于 Paper 01 与 Paper 02 各自关注单一维度（验证信息或约束信息），本研究将这些维度视为一个整体的特征空间，考察其潜在结构。

\subsubsection{2.2 共同基础理论与指令理解}


共同基础理论（Clark \& Marshall, 1981; Clark, 1996; Clark \& Brennan, 1991）指出，成功交际依赖于交际双方共享的背景知识与情境信息。在人机交互中，智能体需要基于用户指令推断当前状态、目标状态以及达到目标所需的操作序列。若指令中明确包含操作对象、操作顺序、初始状态与完成标准，则共同基础更为充分，智能体在规划与验证阶段的不确定性更低。

将这一理论框架与本系列前两项研究对接，可形成以下推论。Paper 01 表明，验证信息为智能体提供了终止判断的基础；Paper 02 表明，约束信息在适度时压缩了动作空间，过度时增加了负荷。共同基础理论则将这些发现整合进一个更为一般的框架：\textbf{指令语用结构的根本功能是为智能体建立"共享任务情境"}。验证信息界定"终点"，约束信息界定"边界"，对象信息界定"实体"，动作信息界定"路径"——各类语用信息的组合方式决定了共同基础的充分程度，进而影响执行路径的稳定性。

\subsubsection{2.3 研究假设}


基于上述理论，本研究提出以下假设：

\textbf{H1（语用类型的存在性）}：AgentNet GUI-Agent 用户指令存在可识别的潜在语用结构类型，可通过 LCA 从 11 个二元语用特征中识别出来。本假设为 LCA 的基本前提——若数据中不存在潜在类别结构，模型选择指标将不支持 $k>1$ 的解。

\textbf{H2（完成率差异）}：不同语用结构类型的任务完成率存在显著差异。若 H2 成立，则说明语用结构的组织方式与任务最终结果存在关联。

\textbf{H3（过程稳定性差异）}：不同语用结构类型在总步数、错误比例、冗余比例以及路径稳定性指标（最长连续错误步数、转换次数、转换率）上存在显著差异。H3 是本研究的核心假设——它关注的是"如何执行"而非"是否完成"，更直接地检验语用结构对执行过程的影响。

\textbf{H4（独立于显式性评分）}：在控制显式性评分（11 个特征之和）后，语用结构类型仍与完成率存在独立关联。若 H4 成立，则说明语用结构的组合方式提供的预测信息无法被简单的特征计数所替代——这是对"组合性"问题的直接检验。

\textbf{H5（任务难度调节效应）}：语用结构类型与任务难度对完成率的关联存在交互效应。基于 Paper 02 的发现——约束密度的效应随任务难度而变化——可预期，包含更多约束信息的语用类型在不同难度下的表现可能存在不对称性。

\subsection{3 研究方法}


\subsubsection{3.1 数据来源与样本}


本研究使用 AgentNet 公开 GUI 交互轨迹数据集（Zhang et al., 2024）。原始数据包含 22,811 条任务尝试记录。经对关键结果变量（`task\_completed`、`total\_steps`、`incorrect\_steps`、`redundant\_ratio`）进行完整案例分析后，最终纳入分析的记录为 20,484 条。分析单元为单次任务尝试；需注意，同一 `task\_id` 可能存在多次尝试，步级路径指标按 `task\_id` 聚合后映射到所有尝试——因此路径分析中有效样本量的解释需审慎。

本研究属于对公开数据的二次分析，未重新运行智能体，亦未收集新的交互日志。

\subsubsection{3.2 语用特征提取}


依据语用类型学与言语行为分类学，本研究从指令文本中提取 11 个二元语用特征，覆盖"谁/什么/怎么做/做到什么程度"四个语用维度（见表 1）。


\begin{table}[htbp]
\centering
\caption{表 1：11 个二元语用特征的定义与语用维度}
\small
\adjustbox{max width=\textwidth}{
\begin{tabular}{l l l}
\toprule
语用维度 & 特征 & 定义 \\
\midrule
谁/什么 & `has\_application` & 是否明确指出应用或软件 \\
什么 & `has\_operation\_object` & 是否明确指出操作对象 \\
什么 & `has\_file\_reference` & 是否提供文件名、路径、工作表名或页面名 \\
什么 & `has\_object\_properties` & 是否给出对象属性（颜色、位置、格式、数量等） \\
什么 & `has\_initial\_state` & 是否说明当前初始状态 \\
怎么做 & `has\_target\_action` & 是否明确指出目标动作 \\
怎么做 & `has\_operation\_order` & 是否说明操作顺序 \\
怎么做 & `has\_multiple\_subtasks` & 是否包含多个可区分的子任务 \\
怎么做 & `has\_constraint\_words` & 是否包含约束词（only / don't / keep / avoid 等） \\
做到什么程度 & `has\_completion\_criteria` & 是否给出完成或验收标准 \\
做到什么程度 & `has\_target\_state` & 是否明确指出目标状态 \\
\bottomrule
\end{tabular}
}
\end{table}

特征提取采用基于关键词与正则表达式的规则方法。显式性评分（`explicitness\_score`）定义为上述 11 个正向特征之和。

关于测量有效性，需说明以下几点。第一，上述特征由规则自动提取，未经人工标注的精度、召回率、F1 或 Cohen's κ 检验——特征命中与否反映的是文本中是否出现对应语言线索，而非对语用功能的完美刻画。第二，部分关键词在不同语境下可能承载不同功能（例如 "check" 既可表示验证标准，也可表示动作本身），规则方法会产生一定误分类。第三，尽管存在上述局限，规则型特征抽取在现有大规模标注数据稀缺的条件下，仍是一种可行且透明的操作化方法。

\subsubsection{3.3 潜类别分析（LCA）实现}


本研究采用潜类别分析（LCA）识别用户指令的潜在语用结构类型。LCA 的核心思想是：观测变量之间的关联可由一个潜在的类别变量解释；在给定潜在类别后，各观测变量条件独立（Lazarsfeld \& Henry, 1968）。

由于 R 语言 `poLCA` 包在当前环境中不可用，本研究使用 Python 实现了标准的二元潜类别 EM 算法（Goodman \& Frank, 2016 的概率推理框架）。对每个 $k=1$ 至 8，分别进行多次随机初始化的 EM 拟合，并保留对数似然最大的解。模型选择综合考察 AIC、BIC、熵（entropy）与 BIC 相对改善率截止准则（改善率低于 0.15\% 时截止）。作为探索性对比，同时拟合了 sklearn 的 GaussianMixture（tied covariance），其在 $k=5$ 时与 EM-LCA 的 Adjusted Rand Index 为 0.372，提示专用二元 LCA 更适合当前数据类型。

EM 算法的实现步骤如下：（1）对每个 $k$，进行多次随机初始化（每次随机分配类别成员概率并归一化）；（2）E 步：基于当前参数估计，计算每条观测属于各类别的后验概率；（3）M 步：基于后验概率加权更新各类别的条件响应概率与类别先验概率；（4）迭代至对数似然收敛（变化量 $< 10^{-6}$）。

\subsubsection{3.4 过程结果与路径指标}


过程结果指标包括：`total\_steps`（总步数）、`incorrect\_ratio`（错误步数占比）、`redundant\_ratio`（冗余步数占比）。路径稳定性指标从步级数据中提取并按 `task\_id` 聚合，包括：`first\_error\_position`（首次错误位置相对总步数）、`max\_consecutive\_incorrect`（最长连续错误步数）、`transition\_count`（相邻正确/错误状态转换次数）与 `transition\_rate`（转换次数占有效相邻对的比例）。

路径稳定性指标的设计遵循以下逻辑：`max\_consecutive\_incorrect` 捕捉智能体偏离正轨的最大持续性；`transition\_count` 与 `transition\_rate` 捕捉正确—错误状态之间的振荡频率——一条稳健的执行轨迹应表现为"少量偏离、快速纠正"，而非"反复错—改—错"的振荡模式。

\subsubsection{3.5 分析流程}


步骤 1：特征提取与描述性统计；
步骤 2：对 $k=1$ 至 8 拟合 EM-LCA，依据 BIC 肘部法则与熵选择最优 $k$；
步骤 3：基于条件响应概率为每个潜在类别赋予语用标签；
步骤 4：使用卡方检验比较各类别完成率，使用 ANOVA/MANOVA 比较过程性结果；
步骤 5：构建控制显式性评分的逻辑回归模型，检验语用类别的独立预测效应；
步骤 6：引入任务难度与语用类别的交互项，检验调节效应；
步骤 7：对执行路径指标进行类别间比较；
步骤 8：通过训练—测试划分与 Bootstrap 进行模型稳定性检查。

\subsection{4 结果}


\subsubsection{4.1 LCA 模型选择}


表 2 报告了 $k=1$ 至 8 的 EM-LCA 拟合统计量。虽然原始 BIC 在 $k=8$ 处达到最小值，但 BIC 相对改善率在 $k=5$ 后急剧下降：$k=5$ → 6 仅改善 0.10\%，$k=6$ → 7 为 0.09\%，$k=7$ → 8 为 0.03\%，均低于 0.15\% 的截止阈值。$k=5$ 时的熵为 0.767，表明类别分离程度可接受。综合 BIC 肘部法则、模型可解释性与研究假设对 3–5 类的预期，最终选择 $k=5$。


\begin{table}[htbp]
\centering
\caption{表 2：LCA 模型选择统计量}
\small
\adjustbox{max width=\textwidth}{
\begin{tabular}{r l r r r r r}
\toprule
$k$ & 对数似然 & AIC & BIC & Entropy & 参数数 & BIC 改善率 \\
\midrule
1 & −110,286.70 & 220,595.41 & 220,682.61 & 1.000 & 11 & — \\
2 & −104,621.78 & 209,289.56 & 209,471.90 & 0.685 & 23 & 5.08\% \\
3 & −101,067.99 & 202,205.97 & 202,483.43 & 0.774 & 35 & 3.34\% \\
4 & −100,205.11 & 200,504.22 & 200,876.81 & 0.787 & 47 & 0.79\% \\
\textbf{5} & \textbf{−99,717.22} & \textbf{199,552.45} & \textbf{200,020.16} & \textbf{0.767} & \textbf{59} & \textbf{0.43\%} \\
6 & −99,554.67 & 199,251.34 & 199,814.19 & 0.802 & 71 & 0.10\% \\
7 & −99,408.79 & 198,983.58 & 199,641.56 & 0.773 & 83 & 0.09\% \\
8 & −99,320.81 & 198,831.62 & 199,584.72 & 0.794 & 95 & 0.03\% \\
\bottomrule
\end{tabular}
}
\end{table}

\begin{figure}[htbp]
\centering
\includegraphics[width=0.92\textwidth]{paper03_figures/fig_bic_elbow.png}
\caption{图 1　BIC 随潜在类别数 $k$ 的变化}
\end{figure}


\textbf{图 1} 直观展示了 BIC 的肘部形态：$k=5$ 后改善率急剧衰减，支持五类别解在拟合优度与简洁性之间达到平衡。H1（语用类型的存在性）得到支持。

\subsubsection{4.2 语用特征剖面与类别标签}


表 3 呈现了五类别的条件响应概率矩阵。图 2 以热图形式直观展示了各类别的特征模式——颜色深浅直接对应条件概率的高低，使类别间的结构差异一目了然。


\begin{table}[htbp]
\centering
\caption{表 3：五类别条件响应概率矩阵}
\small
\adjustbox{max width=\textwidth}{
\begin{tabular}{l r r l r r r r l l l r}
\toprule
类别 & 应用 & 操作对象 & 目标动作 & 目标状态 & 文件引用 & 对象属性 & 操作顺序 & 多子任务 & 约束词 & 完成标准 & 初始状态 \\
\midrule
Object/File-Rich & 0.337 & \textbf{0.973} & 0.850 & 0.020 & \textbf{0.753} & 0.595 & 0.631 & 0.871 & 0.188 & 0.000 & 0.127 \\
Minimal & 0.142 & 0.147 & 0.227 & 0.006 & 0.029 & 0.131 & 0.115 & 0.104 & 0.070 & 0.000 & 0.027 \\
Action-Constrained & 0.251 & 0.381 & \textbf{0.949} & 0.047 & 0.159 & 0.444 & 0.541 & \textbf{0.971} & \textbf{1.000} & \textbf{0.798} & 0.104 \\
Action-Oriented & 0.196 & 0.281 & \textbf{0.834} & 0.031 & 0.069 & 0.371 & 0.441 & \textbf{0.889} & 0.162 & 0.000 & 0.059 \\
Constraint-Verification & 0.100 & 0.092 & 0.330 & 0.021 & 0.022 & 0.166 & 0.310 & 0.557 & \textbf{1.000} & \textbf{0.820} & 0.048 \\
\bottomrule
\end{tabular}
}
\end{table}

\emph{注：} 粗体标示各类别中条件概率 ≥ 0.75 的特征。表中数值为模型估计的概率 $P(\text{feature}=1 \mid \text{class})$。

\begin{figure}[htbp]
\centering
\includegraphics[width=0.92\textwidth]{paper03_figures/fig_profile_heatmap.png}
\caption{图 2　五类别语用特征条件响应概率热图}
\end{figure}


\textbf{图 2} 以颜色梯度呈现了五类别的特征概率模式：行代表类别，列代表特征；颜色越深表示该特征在该类别中的出现概率越高。

基于条件响应概率，本研究将五个潜类别赋予以下语用标签与解释：

\begin{enumerate}
\item[1.] \textbf{Object/File-Rich（对象/文件丰富型）}：操作对象（0.973）与文件引用（0.753）概率最高，操作顺序（0.631）与对象属性（0.595）也较高。该类指令通过明确的实体锚定来降低解析歧义。
\item[2.] \textbf{Minimal（极简型）}：所有特征概率均处于低水平（最高仅 0.227），指令简短且语用信息稀疏。
\item[3.] \textbf{Action-Constrained（动作约束型）}：约束词（1.000）、多子任务（0.971）、目标动作（0.949）与完成标准（0.798）均处于高水平。该类指令同时明确"做什么""做到什么程度"与"不能做什么"。
\item[4.] \textbf{Action-Oriented（动作导向型）}：目标动作（0.834）与多子任务（0.889）概率高，但约束词（0.162）、完成标准（0.000）与文件引用（0.069）概率低。该类指令偏重动作序列的连续性描述。
\item[5.] \textbf{Constraint-Verification（约束验证型）}：约束词（1.000）与完成标准（0.820）概率最高，但目标动作（0.330）与操作对象（0.092）概率低。该类指令侧重边界条件与终点判断，而非路径细节。
\end{enumerate}


\subsubsection{4.3 样本分布与指令示例}


表 4 报告了五类别的实际样本分布。动作导向型占据绝对主导（48.15\%），是最常见的指令类型；约束验证型占比最低（4.19\%）。模型先验概率与观测比例总体一致。


\begin{table}[htbp]
\centering
\caption{表 4：潜在类别样本分布}
\small
\adjustbox{max width=\textwidth}{
\begin{tabular}{l r r}
\toprule
类别 & $n$ & 占比 \\
\midrule
Object/File-Rich（对象/文件丰富型） & 1,631 & 7.96\% \\
Minimal（极简型） & 6,282 & 30.67\% \\
Action-Constrained（动作约束型） & 1,849 & 9.03\% \\
Action-Oriented（动作导向型） & 9,864 & 48.15\% \\
Constraint-Verification（约束验证型） & 858 & 4.19\% \\
\textbf{合计} & \textbf{20,484} & \textbf{100.00\%} \\
\bottomrule
\end{tabular}
}
\end{table}

表 5 列出了各类别的真实指令示例，包括原始英文指令、中文意译、命中特征及完成结果。示例从分析数据中直接抽取，未经人工改写。同一类别内部的指令在完成结果上既有成功也有失败，说明语用结构类型与任务完成之间的关联是概率性而非决定性的。


\begin{table}[htbp]
\centering
\caption{表 5：各类别真实指令示例}
\small
\adjustbox{max width=\textwidth}{
\begin{tabular}{l l l l l}
\toprule
类别 & 原始指令（节选） & 中文意译 & 核心命中特征 & 完成 \\
\midrule
Object/File-Rich & "Please help me open the document '...battery technology' in the Downloads folder with LibreOffice Writer. Change the body text of the 'summary' section to 'Capitalize Every Word'..." & 使用 LibreOffice Writer 打开下载文件夹中的电池技术文档，将摘要部分正文改为每词首字母大写并居中对齐。 & 应用、操作对象、目标动作、文件引用、对象属性、操作顺序、多子任务 & 否 \\
Object/File-Rich & "Open the terminal. First, switch to the 'Desktop/wo' directory. Then, list all the files... Next, rename 'go.java'... to 'gost.java'. Finally, list all the files again for verification." & 打开终端，切换到 Desktop/wo，列出文件，将 go.java 重命名为 gost.java，最后列出文件验证。 & 应用、操作对象、目标动作、文件引用、对象属性、操作顺序、多子任务、初始状态 & 是 \\
Minimal & "Uninstall WhatsApp from MacBook" & 从 MacBook 上卸载 WhatsApp。 & 无 & 是 \\
Minimal & "Can you help me like the latest two videos of my favorite YouTube blogger '李子柒' and comment on 'love u' separately?" & 帮我在 YouTube 上给最喜欢的博主李子柒的最新两条视频点赞并评论。 & 无 & 是 \\
Minimal & "Browse my twitter recommendations and set all the politic related ones uninterested." & 浏览 Twitter 推荐，将与政治相关的内容标记为不感兴趣。 & 无 & 否 \\
Action-Constrained & "I wrote a Python code for drawing flowers, but I named it incorrectly... Open the terminal, go to './Desktop', rename the file 'flower' to 'flower.py', check the directory to confirm if renaming is successful, and then run the code 'flower.py'." & 打开终端，进入 ./Desktop，将文件 flower 重命名为 flower.py，检查目录确认重命名成功，然后运行代码。 & 应用、操作对象、目标动作、目标状态、文件引用、对象属性、操作顺序、多子任务、约束词、完成标准 & 是 \\
Action-Constrained & "I need to use the terminal operation. Please help me move the 2.docx file on the desktop to the wastebasket. Then, please help me check how many files there are in the wastebasket... delete the last file completely." & 使用终端，将桌面 2.docx 移到废纸篓，查看废纸篓文件数，列出所有文件，彻底删除最后一个文件。 & 应用、操作对象、目标动作、目标状态、文件引用、对象属性、操作顺序、多子任务、约束词、完成标准 & 是 \\
Action-Oriented & "Open Uber in Google browser, select Business, select The platform under Overview, select Get started for free, and register..." & 在浏览器中打开 Uber，选择 Business，在 Overview 下选择平台，点击免费开始并注册。 & 目标动作、目标状态、对象属性、操作顺序、多子任务 & 是 \\
Action-Oriented & "In the Google Forms web version, Click on Settings, click to open Make this a quiz... Limit to 1 response..." & 在 Google Forms 网页版中，点击 Settings，开启测验模式，限制 1 次回复。 & 目标动作、目标状态、对象属性、操作顺序、多子任务 & 否 \\
Constraint-Verification & "Check out the official introduction to Elden Ring and check the Steam store to see if it's on sale." & 查看艾尔登法环的官方介绍，检查 Steam 商店是否在打折。 & 约束词、完成标准 & 是 \\
Constraint-Verification & "Disable the uppercase/lowercase check for nouns in IDEA" & 在 IDEA 中禁用名词的大小写检查。 & 约束词、完成标准 & 是 \\
Constraint-Verification & "Change Confirm On Exit in the Debug function of VS Code: to Controls." & 将 VS Code Debug 功能中的 Confirm On Exit 改为 Controls。 & 约束词、完成标准 & 否 \\
\bottomrule
\end{tabular}
}
\end{table}

\subsubsection{4.4 执行路径稳定性差异}


本节聚焦于过程稳定性——本研究最核心的发现维度。表 6 报告了各类别在总步数、错误比例与冗余比例上的均值与标准差；表 7 报告了差异检验结果。


\begin{table}[htbp]
\centering
\caption{表 6：各类别过程性结果指标描述统计}
\small
\adjustbox{max width=\textwidth}{
\begin{tabular}{l l l l}
\toprule
类别 & total\_steps $M$ ($SD$) & incorrect\_ratio $M$ ($SD$) & redundant\_ratio $M$ ($SD$) \\
\midrule
Object/File-Rich & 22.37 (16.33) & 0.121 (0.105) & 0.075 (0.068) \\
Minimal & 18.44 (16.06) & 0.141 (0.128) & 0.093 (0.101) \\
Action-Constrained & 19.34 (9.61) & \textbf{0.110} (0.101) & \textbf{0.064} (0.075) \\
Action-Oriented & 21.70 (17.77) & 0.128 (0.109) & 0.075 (0.079) \\
Constraint-Verification & 20.40 (14.30) & 0.125 (0.120) & 0.085 (0.089) \\
\bottomrule
\end{tabular}
}
\end{table}


\begin{table}[htbp]
\centering
\caption{表 7：过程性结果指标的类别差异检验}
\small
\adjustbox{max width=\textwidth}{
\begin{tabular}{l r r r}
\toprule
变量 & $F$ & ANOVA $p$ & $\eta^2$ \\
\midrule
total\_steps & 45.60 & < 0.001 & 0.009 \\
incorrect\_ratio & 33.32 & < 0.001 & 0.006 \\
redundant\_ratio & 64.53 & < 0.001 & 0.012 \\
MANOVA (Wilks' $\lambda$) & 40.72 & < 0.001 & (1 − $\lambda$) = 0.023 \\
\bottomrule
\end{tabular}
}
\end{table}

MANOVA 结果显示，五类在三个过程指标上的联合差异显著（Wilks' $\lambda = 0.977$，$F(12, 54{,}177) = 40.72$，$p < 0.001$），H3 获得支持。

从描述统计看，一个值得注意的模式浮现：动作约束型在错误比例（0.110）与冗余比例（0.064）上均处于最低水平，而极简型在这两个指标上最高（0.141 和 0.093）。这与"信息越少越简单"的直觉相悖——极简型指令虽然简短，但过程的"效率损失"反而最大，智能体可能需要更多试探性步骤来补全缺失的语用信息。

为进一步刻画执行路径的稳定性，本研究分析了步级聚合指标（表 8、表 9、图 3）。


\begin{table}[htbp]
\centering
\caption{表 8：各类别执行路径指标描述统计}
\small
\adjustbox{max width=\textwidth}{
\begin{tabular}{l l l l l}
\toprule
类别 & first\_error\_position $M$ ($SD$) & max\_consecutive\_incorrect $M$ ($SD$) & transition\_count $M$ ($SD$) & transition\_rate $M$ ($SD$) \\
\midrule
Object/File-Rich & 0.298 (0.254) & 1.236 (0.877) & 4.128 (4.617) & 0.193 (0.133) \\
Minimal & 0.306 (0.284) & 1.262 (1.212) & 3.587 (4.144) & 0.207 (0.159) \\
Action-Constrained & 0.318 (0.290) & \textbf{1.117} (1.001) & \textbf{3.143} (3.115) & \textbf{0.169} (0.136) \\
Action-Oriented & 0.309 (0.278) & 1.275 (1.058) & 4.210 (5.135) & 0.195 (0.140) \\
Constraint-Verification & 0.307 (0.285) & 1.235 (1.185) & 3.826 (4.772) & 0.185 (0.145) \\
\bottomrule
\end{tabular}
}
\end{table}

\emph{注：} `first\_error\_position` 的有效样本量因部分任务无错误步而减少（各类别 $n$ 分别为 1,437 / 5,122 / 1,479 / 8,380 / 691）。


\begin{table}[htbp]
\centering
\caption{表 9：执行路径指标的类别差异检验}
\small
\adjustbox{max width=\textwidth}{
\begin{tabular}{l r r r r}
\toprule
变量 & $F$ & ANOVA $p$ & Kruskal-Wallis $H$ & $p$ \\
\midrule
first\_error\_position & 1.02 & 0.398 & 3.15 & 0.534 \\
max\_consecutive\_incorrect & 8.40 & < 0.001 & 50.27 & < 0.001 \\
transition\_count & 31.55 & < 0.001 & 149.10 & < 0.001 \\
transition\_rate & 25.29 & < 0.001 & 86.43 & < 0.001 \\
\bottomrule
\end{tabular}
}
\end{table}

\begin{figure}[htbp]
\centering
\includegraphics[width=0.92\textwidth]{paper03_figures/fig_path_metrics.png}
\caption{图 3　各类别路径稳定性指标比较}
\end{figure}


\textbf{图 3} 以并排箱线图或条形图展示了各类别在最长连续错误步数、转换次数与转换率上的分布。动作约束型在所有三项指标上均为最低值，呈现出一致的"最稳健"模式。

路径指标揭示了几个值得注意的模式。第一，动作约束型在最长连续错误步数（$M=1.12$）、转换次数（$M=3.14$）与转换率（$M=0.169$）上均最低——综合三项指标，动作约束型的执行轨迹最为稳健，呈现出"少量偏离、快速纠正"的特征。第二，动作导向型的转换次数最高（$M=4.21$），尽管其错误步数并非最高，提示该类指令的执行过程波动较大：智能体重描述"做什么"，但缺少约束与验证信息来指引"何时纠正"。第三，首次错误位置在类别间无显著差异（$F=1.02$，$p=0.398$），说明语用结构类型主要影响的是错误发生后的恢复模式与整体路径波动，而非错误首次出现的早晚——即结构类型的效应作用于"纠错"阶段而非"犯错"阶段。

\subsubsection{4.5 任务完成率比较与显式性独立效应}


表 10 展示了各类别的任务完成率及 95\% 置信区间。约束验证型完成率最高（82.6\%），动作导向型最低（78.3\%）。卡方检验显示各类别完成率差异显著（$\chi^2(4)=26.82$，$p<0.001$），H2 得到支持。


\begin{table}[htbp]
\centering
\caption{表 10：各类别任务完成率及 95\% 置信区间}
\small
\adjustbox{max width=\textwidth}{
\begin{tabular}{l r r r r}
\toprule
类别 & $n$ & 完成数 & 完成率 & 95\% CI \\
\midrule
Object/File-Rich & 1,631 & 1,336 & 81.9\% & [79.97\%, 83.71\%] \\
Minimal & 6,282 & 4,946 & 78.7\% & [77.70\%, 79.73\%] \\
Action-Constrained & 1,849 & 1,513 & 81.8\% & [80.00\%, 83.52\%] \\
Action-Oriented & 9,864 & 7,724 & 78.3\% & [77.48\%, 79.11\%] \\
Constraint-Verification & 858 & 709 & 82.6\% & [79.96\%, 85.02\%] \\
\bottomrule
\end{tabular}
}
\end{table}

需要强调的是，完成率差异虽然显著，但效应量有限——各类别完成率均集中在 78\%–83\% 的窄区间。这一现象本身即传递了一个方法论信号：\textbf{完成率作为一个综合性终局指标，可能无法充分反映语用结构对执行过程的影响}。因此，本研究更关注过程稳定性指标（§4.4）与控制显式性评分后的独立效应（下文）。

为检验语用类别是否提供超出线性显式性评分的信息，本研究拟合了逻辑回归模型：$\text{completed} \sim \text{C(class)} + \text{explicitness\_score}$，以对象/文件丰富型为参照类。表 11 报告了结果。


\begin{table}[htbp]
\centering
\caption{表 11：控制显式性评分后的逻辑回归系数}
\small
\adjustbox{max width=\textwidth}{
\begin{tabular}{l l r l r l}
\toprule
预测变量 & $b$ & SE & $z$ & $p$ & 95\% CI \\
\midrule
Intercept & 1.409 & 0.111 & 12.71 & < 0.001 & [1.192, 1.626] \\
Minimal vs. Object/File-Rich & −0.117 & 0.104 & −1.12 & 0.262 & [−0.320, 0.087] \\
Action-Constrained vs. Object/File-Rich & −0.009 & 0.088 & −0.11 & 0.916 & [−0.182, 0.164] \\
\textbf{Action-Oriented vs. Object/File-Rich} & \textbf{−0.189} & \textbf{0.076} & \textbf{−2.48} & \textbf{0.013} & \textbf{[−0.339, −0.040]} \\
Constraint-Verification vs. Object/File-Rich & 0.089 & 0.116 & 0.77 & 0.443 & [−0.139, 0.317] \\
explicitness\_score & 0.018 & 0.016 & 1.12 & 0.262 & [−0.014, 0.050] \\
\bottomrule
\end{tabular}
}
\end{table}

核心发现：在控制了显式性评分后，类别变量的似然比检验仍然显著（$\chi^2(4)=19.53$，$p<0.001$），H4 得到支持。动作导向型与对象/文件丰富型之间的差异显著（$b=-0.189$，$p=0.013$），而显式性评分本身不显著（$b=0.018$，$p=0.262$）。这一结果具有重要的方法论含义：\textbf{显式性评分作为简单的线性加总指标，不足以捕捉指令语用结构的全部信息}。将 11 个特征简单地相加求和，丢失了特征之间的组合性信息——这正是 LCA 所捕捉的"结构"层面。

\subsubsection{4.6 任务难度的调节效应}


为检验语用类别是否调节任务难度与完成率之间的关联，本研究拟合了交互模型：$\text{completed} \sim \text{C(class)} \times \text{task\_difficulty} + \text{explicitness\_score}$。联合交互似然比检验为边际显著（$\chi^2(4)=8.16$，$p=0.086$）。表 12 报告了交互项系数。


\begin{table}[htbp]
\centering
\caption{表 12：语用类别与任务难度交互模型系数（节选）}
\small
\adjustbox{max width=\textwidth}{
\begin{tabular}{l l r l r l}
\toprule
预测变量 & $b$ & SE & $z$ & $p$ & 95\% CI \\
\midrule
\textbf{Action-Constrained × task\_difficulty} & \textbf{−0.175} & \textbf{0.072} & \textbf{−2.43} & \textbf{0.015} & \textbf{[−0.316, −0.034]} \\
Minimal × task\_difficulty & −0.038 & 0.053 & −0.72 & 0.471 & [−0.142, 0.066] \\
Action-Oriented × task\_difficulty & −0.029 & 0.053 & −0.56 & 0.577 & [−0.132, 0.074] \\
Constraint-Verification × task\_difficulty & −0.073 & 0.079 & −0.93 & 0.354 & [−0.228, 0.081] \\
task\_difficulty & 0.030 & 0.049 & 0.61 & 0.540 & [−0.066, 0.127] \\
explicitness\_score & 0.016 & 0.016 & 0.97 & 0.331 & [−0.016, 0.048] \\
\bottomrule
\end{tabular}
}
\end{table}

\begin{figure}[htbp]
\centering
\includegraphics[width=0.92\textwidth]{paper03_figures/fig_interaction_difficulty.png}
\caption{图 4　任务难度与语用类别对完成率的交互效应}
\end{figure}


\textbf{图 4} 展示了按任务难度预测的各类别完成率曲线。动作约束型（蓝色线）在低难度区间的预测完成率最高，但随难度增加，该优势逐步减弱——与其他类型的差距收窄。

动作约束型与任务难度的交互项显著为负（$b=-0.175$，$p=0.015$）——在五个类别中，仅此一项达到显著水平。H5 得到部分支持。这一发现与 Paper 02 的结论形成了有意义的交叉呼应：Paper 02 发现，高难度任务中对约束密度的容忍度较低（最优约 1.93\%）；Paper 03 则进一步显示，即使语用结构整体上包含丰富的约束与验证信息（动作约束型），其完成率优势也会随难度上升而衰减。两条线索共同指向一个更一般的推论：\textbf{任务复杂度是语用信息效能的关键边界条件}——当任务本身足够复杂，即使指令结构相对完备，执行结果仍受任务固有难度的强力约束。

\subsubsection{4.7 LCA 稳定性检查}


为评估五类别解的可靠性，本研究进行了三项稳定性检查（表 13）。


\begin{table}[htbp]
\centering
\caption{表 13：LCA 稳定性检查汇总}
\small
\adjustbox{max width=\textwidth}{
\begin{tabular}{l l l}
\toprule
检查方法 & 指标 & 结果 \\
\midrule
训练—测试划分（8:2） & Adjusted Rand Index & 0.958 \\
k=4/5/6 比较 & BIC & k=5 为 200,020.21（k=4: 200,877.12; k=6: 199,944.38） \\
Bootstrap（5 次，2 次启动） & ARI 均值 (SD) & 0.859 (0.164, 范围 0.548–0.992) \\
\bottomrule
\end{tabular}
}
\end{table}

\emph{注：} 训练—测试划分：训练集 $n=16{,}388$ 拟合 $k=5$ EM-LCA，利用训练参数预测测试集 $n=4{,}096$ 的后验类别，与全量数据 $k=5$ 解的类别对齐后计算 ARI。Bootstrap 重抽样次数（5 次）与随机启动次数（2 次）均有限，结果仅作初步稳定性参考。

现有稳定性证据支持五类别解具有较好的可重复性。然而，若要达到 publication-grade 标准，建议在充足计算资源下补充：（1）200 次以上 Bootstrap 重抽样，每次至少 30 次随机启动；（2）不同随机种子下的多次完整拟合，报告 ARI 分布；（3）若未来可使用 R/poLCA，补充 Lo-Mendell-Rubin 似然比检验与参数 Bootstrap 置信区间。

\subsection{5 讨论}


\subsubsection{5.1 五种语用结构类型的执行意义}


通过 LCA 识别的五种语用结构类型，可理解为用户在 GUI 任务场景中组织任务信息的五种语用策略。不同于将指令特征视为连续谱系上的任意点，类型学视角揭示了指令语用结构的"离散性"——用户倾向于以特定的特征组合方式来组织指令，而非在所有维度上等概率地"多写一点"。

对象/文件丰富型通过明确的实体锚定降低了实体解析阶段的不确定性。其总步数最长（$M=22.37$），可能反映该类任务本身的复杂度较高（涉及文件操作、对象属性修改等）；但同时，其错误比例与冗余比例处于中等水平——信息丰富并不必然带来高错误率，暗示充分的语用信息可能在更复杂的任务中发挥了"导航"功能。

极简型是第二常用类型（30.67\%），但其过程指标表现最差——错误比例（0.141）与冗余比例（0.093）均为五类最高。这一发现挑战了"指令越短越好"的朴素直觉：极简型指令虽然节省了用户的输入成本，但将信息补全的认知成本转移给了智能体。在缺少明确操作对象、动作顺序与验证标准的情况下，智能体需要通过试探—纠错来逼近目标，表现为较高的冗余步比例与路径转换率。

动作约束型在过程稳定性上的一致优势（三项路径指标均为最低）是本研究最值得关注的发现。该类指令同时具备明确的目标动作（0.949）、完成标准（0.798）与约束词（1.000）——换言之，它同时在"做什么""做到什么程度"与"不能做什么"三个语用维度上为智能体提供了信息。这一"三维完整"的语用结构可能使智能体在规划、执行与验证三个阶段均获得明确指引，从而减少了执行过程中的方向性偏差与振荡。

动作导向型是绝对主导的指令类型（48.15\%），但其完成率最低、转换次数最高。这一模式提示，"动作描述丰富"并不等同于"语用结构完整"——如果大量的动作描述缺乏约束边界与验证终点的配合，智能体可能在执行过程中更容易产生路径波动。

约束验证型样本量最小（4.19\%）但完成率最高（82.6\%）。该类指令通常对应"检查""确认""设置变更"等目标状态较易判定的任务。然而，受限于样本量与任务类型分布，对其过程优势仍需谨慎解读。

\subsubsection{5.2 路径稳定性作为语用结构效应的核心证据}


本研究主张，评估指令质量应延伸至执行过程的稳定性，而非止步于完成率。完成率是综合性终局指标，同时受任务难度、智能体能力与环境状态等多种因素影响；相较之下，路径稳定性指标——最长连续错误步数、转换次数与转换率——更为直接地反映了指令语用结构对执行过程的微观影响。

本研究中，动作约束型在完成率上并非最高（81.8\%，低于约束验证型的 82.6\%），但其路径稳定性最优。这一"完成率与路径稳定性脱耦"现象提示：不同的语用结构可能通过不同的机制影响执行效果——约束验证型的完成率高可能得益于任务本身目标状态易于判定；动作约束型的路径稳定性好则可能源于执行过程中的方向性指引更充分。仅关注"是否完成"会掩盖语用结构在"如何执行"层面的系统性差异。

\subsubsection{5.3 与 Paper 01/02 的理论整合}


将本研究的发现与前两项研究整合，可构建一个更为完整的理论图景（图 5）。

Paper 01 表明，验证信息（完成标准/验证条件）与执行效率正向关联——验证信息为智能体提供了终止判准。Paper 02 表明，约束信息的效应是条件性的——在高难度或长指令任务中呈现倒 U 形，最优密度约 1.7\%–1.9\%。Paper 03 则揭示了更高层次的组织原则：用户指令中的不同语用维度不是随机组合的，而是倾向于聚合成若干可识别的"语用结构类型"——在这些类型中，动作导向型（以"怎么做"为主导但缺乏约束与验证）完成率最低，动作约束型（同时包含"怎么做""做到什么程度"与"不能做什么"）路径最稳健。

三者共同指向以下推论：\textbf{用户指令对智能体执行的影响，既取决于关键语用成分的剂量（密度），也取决于这些成分的组织方式（结构类型），最终都受任务复杂度的调节。} 这一推论具有直接的实践含义——提示工程不应仅关注"加了什么信息"，还应关注"信息如何组织"。

\begin{figure}[htbp]
\centering
\includegraphics[width=0.92\textwidth]{paper03_figures/fig5_conceptual_framework.png}
\caption{图 5　三篇论文的理论整合框架}
\end{figure}


\textbf{图 5} 以概念框架图的形式整合了本系列三项研究的理论关系。自下而上呈现了三个递进的分析层次——层次一（单一语用成分的存在/缺失：Paper 01 验证信息与 Paper 02 Stage 1 约束信息）、层次二（语用成分的剂量效应：Paper 02 约束密度的非线性与条件效应）、层次三（整体语用结构类型：Paper 03 五类别 LCA）。右侧标注了贯穿三个层次的情境调节变量（任务复杂度），底部汇总了系列研究所覆盖的多维度执行效果（完成率、执行步数、路径稳定性、冗余步数比、错误模式）。该图揭示的核心推论是：指令效果 = f（语用成分剂量，语用结构类型，任务复杂度）。

\subsubsection{5.4 显式性评分的局限性}


控制显式性评分后，语用类别仍与完成率存在独立关联（$\chi^2(4)=19.53$，$p<0.001$），而显式性评分本身不显著（$b=0.018$，$p=0.262$）。这一结果为"组合性"问题提供了经验证据：\textbf{不同语用维度的组合方式提供的预测信息，无法被各维度的简单加总所替代。} 举例而言，一条得分为 5 的指令可能对应完全不同的语用结构——可能是"5 个动作描述"（动作导向型），也可能是"对象+动作+顺序+约束+验证"（动作约束型）——两者在执行效果上可能存在系统性差异。

这一发现对提示工程的启示是：优化提示时不宜简单追求"显式性分数"的提高，而应关注语用结构的完整性——确保指令在"谁/什么""怎么做""做到什么程度""不能做什么"四个维度上均有覆盖。

\subsubsection{5.5 局限与未来方向}


本研究存在以下局限：

\begin{enumerate}
\item[1.] \textbf{规则型特征的有效性}：11 个二元语用特征由关键词与正则表达式规则提取，未经人工标注的精度、召回率、F1 或 Cohen's κ 检验。未来可引入大语言模型辅助标注或进行人工校验，提升特征有效性。

\item[2.] \textbf{观测数据的因果解释}：语用结构类型与执行效果之间的关联不能解释为因果效应。用户可能根据任务复杂度自选择不同的表达方式，类别—结果关联可能受任务特征混淆。未来可通过实验操纵（如随机分配不同语用结构的指令）或工具变量方法进一步探讨因果路径。

\item[3.] \textbf{LCA 实现的近似性}：自定义 EM 算法未提供 Lo-Mendell-Rubin 检验与参数 Bootstrap 置信区间；稳定性检查虽初步支持五类别解，但 Bootstrap 次数有限（5 次），publication-grade 推断需更充分的计算验证。

\item[4.] \textbf{结论的外部效度}：本研究基于 AgentNet 桌面 GUI 轨迹数据，向网页智能体、移动设备、机器人控制等环境的推广需进一步验证。

\item[5.] \textbf{重复 task\_id 与路径指标}：同一 `task\_id` 的多次尝试导致路径分析的有效样本量解释需审慎，未来可采用混合效应模型处理任务层级嵌套结构。
\end{enumerate}


\subsection{6 结论与展望}


本研究以 AgentNet 大规模 GUI 轨迹数据为基础，从"语用结构类型学"的视角考察了用户指令的语用组织方式及其执行效应。核心发现如下：

\begin{enumerate}
\item[1.] 任务型交互智能体的用户指令存在五种可识别的潜在语用结构类型：对象/文件丰富型、极简型、动作约束型、动作导向型与约束验证型。
\item[2.] 五类在执行路径稳定性、过程性结果与任务完成率上均存在显著差异。动作约束型在执行路径上最为稳健——最长连续错误步数、转换次数与转换率均为最低；极简型的过程效率损失最大。
\item[3.] 语用结构类型对完成率的预测关系独立于线性显式性评分，表明特征的组合方式比特征的简单计数包含更丰富的预测信息。
\item[4.] 动作约束型在低难度任务中的完成率优势随难度上升而减弱，提示任务复杂度是语用结构效能的关键边界条件。
\end{enumerate}


本研究的理论贡献在于将指令质量评估从"单一成分剂量"（Paper 01/02 的视角）提升至"整体结构类型"的分析层面，并主张以路径稳定性而非仅完成率作为语用结构效应的核心判断指标。实践贡献在于为提示工程提供了基于类型学的诊断框架——不仅看指令"写了多少"，更看指令"如何组织"。

未来研究可从以下方向拓展：第一，将语用结构类型作为先验信息融入智能体的自适应提示补全模块，实现面向语用结构的动态指令优化；第二，在跨语言、跨平台数据中检验该类型学的跨域稳健性；第三，采用实验设计检验语用结构类型与执行效果之间的因果路径；第四，结合本系列 Paper 01（验证密度）与 Paper 02（约束密度）的发现，构建一个整合"成分剂量 × 结构类型 × 任务复杂度"的多层次指令效能预测框架。


\newpage

\chapter{综合讨论与展望}


\section{核心发现回顾}


本系列三项研究构建了一个递进式的分析框架，分别从"单一语用成分的存在性"、"语用成分的密度/剂量效应"和"整体语用结构类型"三个层次，考察了用户指令语用特征对任务型交互智能体执行效果的影响。

研究一表明，验证信息（完成标准）作为终止判准，提供了智能体判断"何时停止"的关键语言线索。研究二进一步揭示，约束信息的效应并非简单的线性关系，而是取决于任务情境——在高难度和长指令条件下呈现显著倒 U 形。研究三将分析单元提升至语用结构的整体组织方式，发现了五类稳定的指令类型，其中动作约束型在路径稳定性上表现最优。

\section{理论贡献}


\begin{enumerate}
\item[1.] \textbf{言语行为理论的拓展}：将 Austin 和 Searle 的经典言语行为理论拓展至人机交互领域，揭示了指令内部语用结构的多样性及其行为后果。
\item[2.] \textbf{关联理论的量化检验}：为 Sperber 和 Wilson 的关联理论在人机交互情境中的应用提供了大规模经验证据。验证信息降低认知成本、约束信息存在"过犹不及"效应——两者从不同侧面支撑了认知经济性原则。
\item[3.] \textbf{认知负荷理论在智能体系统中的延伸}：证明了即便对于 LLM 驱动的智能体系统，输入端的语用密度仍会通过规划、验证和冲突消解机制影响任务执行效果。
\item[4.] \textbf{共同基础理论的过程层面实证}：Clark 和 Brennan 的共同基础理论被拓展至智能体执行的过程层面——语用结构的充分性不仅影响"能否完成"，更影响"如何完成"（路径稳定性）。
\end{enumerate}


\section{实践启示}


\begin{enumerate}
\item[1.] \textbf{指令设计原则}：对高难度、长指令任务，约束密度建议控制在 1.7\%–1.9\% 左右，避免约束过载。
\item[2.] \textbf{完成标准的价值}：明确包含验证条件的指令与更高的执行效率相关联，建议在提示工程和指令模板中嵌入明确的完成判据。
\item[3.] \textbf{语用结构诊断}：五类语用结构类型为用户指令质量评估提供了类型学框架，可用于智能体提示优化的前置诊断。
\item[4.] \textbf{终止机制设计}：智能体系统应增强对指令中完成标准的识别与利用能力，将终止判断从启发式规则升级为语言驱动的内部终止规则。
\end{enumerate}


\section{局限与未来方向}


\begin{enumerate}
\item[1.] \textbf{观测性数据局限}：AgentNet 为观测性轨迹数据，所有发现均为统计关联而非因果效应。未来研究可采用实验设计，随机分配约束密度和完成标准条件，实现因果推断。
\item[2.] \textbf{数据集代表性}：AgentNet 的指令主要来源于英文环境下的桌面/网页 GUI 任务。中文及跨语言、跨操作系统场景的适用性有待验证。
\item[3.] \textbf{智能体架构差异}：不同智能体系统（如基于视觉的 vs. 基于 DOM 的）可能对同一语用特征产生不同响应，本研究未能区分架构类型的调节效应。
\item[4.] \textbf{语用特征细粒度}：11 个二元语用特征虽覆盖了主要维度，但未能捕捉语义强度、约束优先级、歧义程度等连续谱系信息。
\item[5.] \textbf{动态适应}：约束密度的最优值可能随智能体能力演进和任务域迁移而变化，未来需建立自适应约束密度估计框架。
\item[6.] \textbf{终止机制的因果检验}：本研究提出的"完成标准→终止判准→执行效率"理论链尚需通过实验操控来检验其因果方向。
\end{enumerate}


\newpage

\begin{thebibliography}{99}


\bibitem{austin1962} Austin, J. L. (1962). \emph{How to Do Things with Words}. Oxford University Press.

\bibitem{clark1991} Clark, H. H., \& Brennan, S. E. (1991). Grounding in communication. In \emph{Perspectives on socially shared cognition} (pp. 127–149).

\bibitem{gigerenzer1999} Gigerenzer, G., Todd, P. M., \& the ABC Research Group. (1999). \emph{Simple Heuristics That Make Us Smart}. Oxford University Press.

\bibitem{goedert2013} Goedert, M., et al. (2013). Instruction clarity and performance in human-computer interaction. \emph{Human Factors, 55}(4), 772–783.

\bibitem{lasecki2013} Lasecki, W. S., et al. (2013). Crowdsourcing step-by-step instructions for GUI tasks. In \emph{Proceedings of UIST}.

\bibitem{liang2020} Liang, J., et al. (2020). Language models are few-shot learners for GUI tasks. In \emph{Proceedings of ACL}.

\bibitem{li2024} Li, X., et al. (2024). A survey of GUI agents: Methods, applications, and challenges. \emph{arXiv preprint}.

\bibitem{pal2022} Pal, S., et al. (2022). WebGPT: Browser-assisted question-answering with human feedback. In \emph{Advances in Neural Information Processing Systems, 35}.

\bibitem{payne1993} Payne, J. W., Bettman, J. R., \& Johnson, E. J. (1993). \emph{The Adaptive Decision Maker}. Cambridge University Press.

\bibitem{schatz2003} Schatz, B. R., et al. (2003). Task complexity and instruction design in human-machine interaction. \emph{International Journal of Human-Computer Studies}.

\bibitem{searle1969} Searle, J. R. (1969). \emph{Speech Acts: An Essay in the Philosophy of Language}. Cambridge University Press.

\bibitem{sperber1995} Sperber, D., \& Wilson, D. (1995). \emph{Relevance: Communication and Cognition} (2nd ed.). Blackwell.

\bibitem{yang2022} Yang, Y., et al. (2022). Mind2Web: Towards a generalist agent for the web. In \emph{Advances in Neural Information Processing Systems, 35}.

\bibitem{yuan2023} Yuan, X., et al. (2023). AUSTIN: Training autonomous GUI agents with visual understanding. In \emph{Proceedings of CVPR}.

\bibitem{zhang2024} Zhang, Y., et al. (2024). AgentNet: A large-scale dataset of GUI agent trajectories. \emph{arXiv preprint}.


\end{thebibliography}

\section{补充数据文件}


\begin{itemize}
\item `revised/paper01\_tables/terminal\_phase\_metrics\_by\_task.csv`：每条任务的终止阶段指标。

\item `revised/paper01\_tables/terminal\_phase\_summary\_by\_criteria.csv`：按是否含完成标准分组的描述性汇总。

\item `revised/paper01\_tables/terminal\_phase\_group\_comparisons.csv`：终止阶段指标的组间比较统计量。
\end{itemize}


Dawson, J. F. (2014). Moderation in management research: What, why, when, and how. \emph{Journal of Business and Psychology}, \emph{29}(1), 1–19.

Grice, H. P. (1975). Logic and conversation. In \emph{Speech acts} (pp. 41–58). Brill.

Sweller, J. (1988). Cognitive load during problem solving: Effects on learning. \emph{Cognitive Science}, \emph{12}(2), 257–285.

Yeung, D. Y. (2016). Too much of a good thing? The effectiveness of self-regulated learning interventions in workplace training. \emph{Human Resource Development Review}, \emph{15}(2), 143–168.

\section{附录：补充表格}


\subsection{表 8　Stage 2 约束任务子样本的完整系数表}


\begin{table}[htbp]
\centering
\caption{表 8：Stage 2 约束任务子样本的完整系数表}
\small
\adjustbox{max width=\textwidth}{
\begin{tabular}{l l r r r l r}
\toprule
变量 & Coef & SE & OR & 95\% CI OR & $z$ & $p$ \\
\midrule
Intercept & 3.163 & 0.650 & 23.645 & [6.613, 84.546] & 4.866 & < 0.001 \\
os [win\_mac] & −1.057 & 0.254 & 0.347 & [0.211, 0.572] & −4.162 & < 0.001 \\
constraint\_density & −0.267 & 0.126 & 0.765 & [0.598, 0.980] & −2.123 & 0.034 \\
I(constraint\_density²) & 0.016 & 0.009 & 1.016 & [0.998, 1.034] & 1.735 & 0.083 \\
task\_difficulty & −0.015 & 0.064 & 0.985 & [0.869, 1.116] & −0.238 & 0.812 \\
word\_count & −0.004 & 0.005 & 0.996 & [0.987, 1.006] & −0.737 & 0.461 \\
\bottomrule
\end{tabular}
}
\end{table}

\subsection{表 9　Johnson-Neyman 风格的简单斜率与曲率（交互模型）}


\begin{table}[htbp]
\centering
\caption{表 9：Johnson-Neyman 风格的简单斜率与曲率（交互模型）}
\small
\adjustbox{max width=\textwidth}{
\begin{tabular}{r l r l r l r l r}
\toprule
任务难度 & 密度=0 处斜率 & SE & $z$ & $p$ & 曲率 $b_2 + b_{2m}\cdot D$ & SE & $z$ & $p$ \\
\midrule
0 & −0.167 & 0.157 & −1.062 & 0.288 & 0.036 & 0.020 & 1.849 & 0.064 \\
1 & −0.131 & 0.132 & −0.994 & 0.320 & 0.027 & 0.016 & 1.692 & 0.091 \\
2 & −0.095 & 0.107 & −0.887 & 0.375 & 0.018 & 0.013 & 1.423 & 0.155 \\
3 & −0.059 & 0.084 & −0.705 & 0.481 & 0.009 & 0.010 & 0.923 & 0.356 \\
4 & −0.023 & 0.063 & −0.365 & 0.715 & −0.000 & 0.008 & −0.023 & 0.982 \\
5 & 0.013 & 0.049 & 0.268 & 0.789 & −0.009 & 0.007 & −1.319 & 0.187 \\
6 & 0.049 & 0.047 & 1.039 & 0.299 & −0.018 & 0.009 & −2.162 & 0.031 \\
7 & 0.085 & 0.059 & 1.432 & 0.152 & −0.028 & 0.011 & −2.459 & 0.014 \\
8 & 0.121 & 0.079 & 1.530 & 0.126 & −0.037 & 0.014 & −2.540 & 0.011 \\
9 & 0.157 & 0.102 & 1.538 & 0.124 & −0.046 & 0.018 & −2.555 & 0.011 \\
10 & 0.193 & 0.127 & 1.524 & 0.128 & −0.055 & 0.022 & −2.549 & 0.011 \\
\bottomrule
\end{tabular}
}
\end{table}

\emph{注：} 表中为沿任务难度取值范围（0–10）选取的整数值。曲率显著为负（$p<0.05$）的区域大致对应任务难度 ≥ 6，与分样本结果一致。

\subsection{表 10　探索性过程结果模型}


\begin{table}[htbp]
\centering
\caption{表 10：探索性过程结果模型}
\small
\adjustbox{max width=\textwidth}{
\begin{tabular}{l l l r l r l}
\toprule
模型 & 结果变量 & 约束密度 $b$ & $p$ & 约束密度² $b_2$ & $p$ & 备注 \\
\midrule
负二项回归 & total\_steps & −0.013 & 0.065 & 0.00010 & 0.915 & 不显著 \\
OLS & redundant\_ratio & 0.002 & 0.127 & −0.00038 & 0.012 & 显著倒 U 形 \\
\bottomrule
\end{tabular}
}
\end{table}

Clark, H. H. (1996). \emph{Using Language}. Cambridge University Press.

Clark, H. H., \& Marshall, C. R. (1981). Definite reference and mutual knowledge. In A. K. Joshi, B. L. Webber, \& I. A. Sag (Eds.), \emph{Elements of Discourse Understanding} (pp. 10–63). Cambridge University Press.

Goodman, N. D., \& Frank, M. C. (2016). Pragmatic language interpretation as probabilistic inference. \emph{Trends in Cognitive Sciences}, \emph{20}(11), 818–829.

Lazarsfeld, P. F., \& Henry, N. W. (1968). \emph{Latent Structure Analysis}. Houghton Mifflin.

Searle, J. R. (1976). A classification of illocutionary acts. \emph{Language in Society}, \emph{5}(1), 1–23.

\section{附录：补充表格与说明}


\subsection{表 A1　k=4/5/6 拟合统计量比较}


\begin{table}[htbp]
\centering
\caption{表 A1：k=4/5/6 拟合统计量比较}
\small
\adjustbox{max width=\textwidth}{
\begin{tabular}{r l r r r r}
\toprule
$k$ & 对数似然 & AIC & BIC & Entropy & 参数数 \\
\midrule
4 & −100,205.27 & 200,504.53 & 200,877.12 & 0.781 & 47 \\
\textbf{5} & \textbf{−99,717.25} & \textbf{199,552.50} & \textbf{200,020.21} & \textbf{0.767} & \textbf{59} \\
6 & −99,619.77 & 199,381.53 & 199,944.38 & 0.759 & 71 \\
\bottomrule
\end{tabular}
}
\end{table}

\subsection{表 A2　补充图表文件索引}


本论文使用的图表保存于项目目录下：

\begin{itemize}
\item `paper03\_figures/fig\_bic\_elbow.png` — BIC 肘部法则图

\item `paper03\_figures/fig\_profile\_heatmap.png` — 语用特征剖面热图

\item `paper03\_figures/fig\_path\_metrics.png` — 路径稳定性指标比较图

\item `paper03\_figures/fig\_interaction\_difficulty.png` — 任务难度交互效应图

\item `paper03\_tables/tab\_model\_selection.csv` — 模型选择统计量

\item `paper03\_tables/tab\_class\_profiles.csv` — 条件响应概率矩阵

\item `paper03\_tables/tab\_examples.json` — 指令示例

\item `paper03\_tables/tab\_stability\_checks.json` — 稳定性检查结果
\end{itemize}


\newpage

\appendix
\chapter{附录}


\section{附录 A　Paper 01 补充数据文件索引}


\begin{itemize}
\item `paper01\_tables/terminal\_phase\_metrics\_by\_task.csv`：每条任务的终止阶段指标
\item `paper01\_tables/terminal\_phase\_summary\_by\_criteria.csv`：按是否含完成标准分组的描述性汇总
\item `paper01\_tables/terminal\_phase\_group\_comparisons.csv`：终止阶段指标的组间比较统计量
\end{itemize}


\section{附录 B　Paper 02 补充数据文件索引}


\begin{itemize}
\item `paper02\_tables/table1\_constraint\_taxonomy.md` — 约束分类法
\item `paper02\_tables/table4\_stage1\_stage2\_models.csv` — Stage 1/2 模型结果
\item `paper02\_tables/table5\_conditional\_effects.csv` — 条件效应模型
\item `paper02\_tables/table8\_stage2\_full\_coefficients.csv` — Stage 2 完整系数
\item `paper02\_tables/table9\_johnson\_neyman\_simple\_slopes.csv` — Johnson-Neyman 简单斜率
\end{itemize}


\section{附录 C　Paper 03 补充数据文件索引}


\begin{itemize}
\item `paper03\_tables/tab\_model\_selection.csv` — LCA 模型选择统计量
\item `paper03\_tables/tab\_class\_profiles.csv` — 类别条件响应概率
\item `paper03\_tables/tab\_examples.json` — 各类别代表性指令示例
\item `paper03\_tables/tab\_stability\_checks.json` — LCA 稳定性检查结果
\end{itemize}


\end{document}